%%% Settings %%%%%%%%%%%%%%%%%%%%%%%%%%%%%%%%%%%%%%%%%%%%%%%%%%%%%%%%%%%%%%%%%%%%%%%%%
\documentclass{article}

\usepackage{graphicx}  % serve per inserire immagini
\usepackage{fancyhdr}  % creazione header-footer
\usepackage{tabularx}  % serve per creare tabelle con colonne a larghezza variabile
\usepackage{ifthen}  % serve per mostrare cose diverse in base a condizioni
\usepackage{geometry}
\usepackage{setspace}
\usepackage{tikz}
\usepackage[italian]{babel}
\usepackage[hidelinks]{hyperref}
\usepackage{pgfgantt}  % per i diagrammi di Gantt
\usepackage{eurosym}
\usepackage{float}
\usepackage{pgf-pie}
\usepackage{pgfplots}
\usepackage{tikz}
\usepackage{pgfplotstable}

% definizione dei colori per i grafici
\definecolor{CO}{HTML}{93b7be}  
\definecolor{DIP}{HTML}{90be6d}  
\definecolor{DIV}{HTML}{6BAED6}  
\definecolor{LA}{HTML}{ffd373}  
\definecolor{PA}{HTML}{a874d2}  
\definecolor{PI}{HTML}{FF8F00}  
\definecolor{RI}{HTML}{a98467}  
\definecolor{extra}{HTML}{E53946}

% setta a 1 se il verbale è esterno, 0 se è interno
\newcommand{\isEsterno}{1}

% Margini della pagina
\geometry{a4paper, margin=1in}

% Intestazione personalizzata
\pagestyle{fancy}
\fancyhf{}
\fancyhead[L]{Code7Crusaders - Software Development Team}
\fancyhead[R]{\thepage}

% Spaziatura delle righe
\setstretch{1.2}

\begin{document}
%%%%%%%%%%%%%%%%%%%%%%%%%%%%%%%%%%%%%%%%%%%%%%%%%%%%%%%%%%%%%%%%%%%%%%%%%%%%%%%%%%%%%%



%%% Sezione del titolo %%%%%%%%%%%%%%%%%%%%%%%%%%%%%%%%%%%%%%%%%%%%%%%%%%%%%%%%%%%%%%%
\begin{titlepage}

    \AddToHookNext{shipout/background}{
        \begin{tikzpicture}[remember picture,overlay]
        \node at (current page.center) {
            \includegraphics{../../../img/background.png}
        };
        \end{tikzpicture}
    }

    \centering
    \vspace*{2cm}

    \includegraphics[width=0.3\textwidth]{../../../img/logo/7Crusaders_logo.png} % logo
    \vspace{1cm}

    {\Huge \textbf{Code7Crusaders}}\\
    \vspace{0.5cm}
    {\Large Software Development Team}\\
    \vspace{2cm}

    {\large \href{https://code7crusaders.github.io/docs/PB/documentazione_interna/glossario.html#piano-di-progetto}{\textbf{Piano di Progetto}\textsuperscript{G}}}\\
    \vspace{5cm}


    \textbf{Membri del Team:}\\
    Enrico Cotti Cottini, Gabriele Di Pietro, Tommaso Diviesti \\
    Francesco Lapenna, Matthew Pan, Eddy Pinarello, Filippo Rizzolo \\
    \vspace{0.5cm}

    \vspace{1cm}
\end{titlepage}
%%%%%%%%%%%%%%%%%%%%%%%%%%%%%%%%%%%%%%%%%%%%%%%%%%%%%%%%%%%%%%%%%%%%%%%%%%%%%%%%%%%%%%



% Versioni %%%%%%%%%%%%%%%%%%%%%%%%%%%%%%%%%%%%%%%%%%%%%%%%%%%%%%%%%%%%%%%%%%%%%%%%%%%
\newpage
\begin{table}[h!]
\centering
\textbf{Versioni} \\ % Titolo sopra la tabella
\vspace{2mm} % Spazio tra il titolo e la tabella
\renewcommand{\arraystretch}{1.7}
    
\begin{tabular}{|c|c|c|c|m{0.35\textwidth}|}
    \hline
    \textbf{Ver.} & \textbf{Data} & \textbf{Autore} & \textbf{Verificatore} & \textbf{Descrizione} \\
    \hline
    1.4 & 14/03/2025 & Lapenna Francesco &  & Compilazione Sprint \\
    \hline
    1.3 & 14/03/2025 & Diviesti Tommaso & Filippo Rizzolo & Aggiunta dei grafici di preventivo e consuntivo agli Sprint \\
    \hline
    1.2 & 13/03/2025 & Diviesti Tommaso & Filippo Rizzolo & Completamento analisi dei rischi: attualizzazione di riscontro rispetto alla loro occorrenza effettiva per i vari Sprint \\
    \hline
    1.1 & 10/03/2025 & Diviesti Tommaso & Filippo Rizzolo & Inizio correzioni segnalate post RTB: sezione "Retrospettiva di periodo (RTB)" \\
    \hline
    1.0 & 19/02/2025 & Diviesti Tommaso & Francesco Lapenna & Approvazione documento \\
    \hline
    0.8 & 17/02/2025 & Francesco Lapenna & Enrico Cotti Cottini & Compilazione sprint \\
    \hline
    0.7 & 5/02/2025 & Francesco Lapenna & Eddy Pinarello & Metriche rilevamento rischi \\
    \hline
    0.6 & 18/01/2025 & Francesco Lapenna & Eddy Pinarello & Rimozione sezione 2.5 \\
    \hline
    0.5 & 18/01/2025 & Diviesti Tommaso & Lapenna Francesco & Aggiornamento analisi rischi \\
    \hline
    0.4 & 7/01/2025 & Diviesti Tommaso & Lapenna Francesco & Modello, vantaggi e caratteristiche sprint \\
    \hline
    0.3 & 27/12/2024 & Lapenna Francesco & Matthew Pan & Inizio compilazione sprint \\
    \hline
    0.2 & 3/12/2024 & Diviesti Tommaso & Lapenna Francesco & Continuazione e revisione del documento \\
    \hline
    0.1 & 29/11/2024 & Lapenna Francesco & Diviesti Tommaso & Prima stesura del documento \\
    \hline
\end{tabular}
\end{table}
%%%%%%%%%%%%%%%%%%%%%%%%%%%%%%%%%%%%%%%%%%%%%%%%%%%%%%%%%%%%%%%%%%%%%%%%%%%%%%%%%%%%%%



% Indice %%%%%%%%%%%%%%%%%%%%%%%%%%%%%%%%%%%%%%%%%%%%%%%%%%%%%%%%%%%%%%%%%%%%%%%%%%%%%
\newpage
\tableofcontents
\listoftables
\listoffigures
%%%%%%%%%%%%%%%%%%%%%%%%%%%%%%%%%%%%%%%%%%%%%%%%%%%%%%%%%%%%%%%%%%%%%%%%%%%%%%%%%%%%%%



% Sezione Introduzione %%%%%%%%%%%%%%%%%%%%%%%%%%%%%%%%%%%%%%%%%%%%%%%%%%%%%%%%%%%%%%%
\newpage
\section{Introduzione}

    \subsection{Scopo del documento}
    Questo documento ha lo scopo di fornire una guida dettagliata e strutturata su come
    il progetto verrà eseguito e gestito. In particolare, verrano trattati i seguenti
    \begin{itemize}
        \item Analisi del Capitolato
        \item Pianificazione
        \item Analisi dei Rischi
        \item Gestione del modello di sviluppo
    \end{itemize}

    \subsection{Scopo del prodotto}
    Il prodotto consiste in una webapp avanzata che integra una chatbot alimentata da intelligenza artificiale,
    pensata per fornire informazioni precise e approfondite su una vasta selezione di prodotti. L’obiettivo principale è
    offrire alle aziende uno strumento semplice ed efficace per accedere a dettagli fondamentali riguardo i prodotti che
    desiderano acquistare, assicurando maggiore trasparenza e chiarezza in ogni fase del processo di selezione. Grazie alla
    nostra soluzione, le aziende possono ottenere risposte immediate su una serie di parametri chiave e informazioni su bibite e relativi produttori/venditori.
    Tutto ciò permette di ridurre incertezze e ambiguità, riducendo i rischi delle aziende legati alla scelta di prodotti non adatti alle proprie esigenze.
    Inoltre, le chat recenti vengono salvate e rese facilmente accessibili agli utenti, permettendo loro di rivedere in qualsiasi
    momento le informazioni precedentemente richieste. Questa funzionalità risulta particolarmente utile per consultare rapidamente
    risposte a domande frequenti o per confrontare dettagli su diversi prodotti, senza dover rifare ogni ricerca, garantendo così
    un'esperienza più efficiente e personalizzata.

    \subsection{Glossario}
    Per avere maggiore chiarezza ed evitare ambiguità per quanto riguarda i termini utilizzati all'interno dei vari documenti,
    viene adottato un \href{https://code7crusaders.github.io/docs/PB/documentazione_interna/glossario.html#glossario}{Glossario\textsuperscript{G}} che contiene una serie di termini e relativa definizione.
    Grazie ad esso, sarà possibile cliccare su una determinata porzione di testo, evidenziata grazie ad uno stile specifico,
    all'interno di un qualsiasi documento e in questo modo si potrà visualizzare la sua definizione all'interno
    del \href{https://code7crusaders.github.io/docs/PB/documentazione_interna/glossario.html#glossario}{Glossario\textsuperscript{G}}

    \subsection{Riferimenti}
    % TODO

    \subsection{Preventivo iniziale}
    Il preventivo iniziale è stato presentato durante la fase di Candidatura ed è pari a \textbf{12805}\euro.
    \\ Per ulteriori informazioni è possibile visualizzare il documento di analisi dei costi e assunzione impegni al seguente link:
    \\ \url{https://code7crusaders.github.io/docs/Candidatura/Preventivo_costi.html}

%%%%%%%%%%%%%%%%%%%%%%%%%%%%%%%%%%%%%%%%%%%%%%%%%%%%%%%%%%%%%%%%%%%%%%%%%%%%%%%%%%%%%%



% Analisi del Capitolato %%%%%%%%%%%%%%%%%%%%%%%%%%%%%%%%%%%%%%%%%%%%%%%%%%%%%%%%%%%%%
\newpage
\section{Analisi del Capitolato}
    \subsection{Obbiettivi del progetto}
    \begin{itemize}
        \item Realizzare un Assistente Virtuale che supporti i clienti nella ricerca
        di informazioni sui prodotti disponibili in catalogo.
        \item Automatizzare le risposte alle domande più frequenti, migliorando
        l'efficienza del servizio clienti.
        \item Integrare un modello \href{https://code7crusaders.github.io/docs/PB/documentazione_interna/glossario.html#llm-large-language-model}{LLM\textsuperscript{G}} esistente per garantire risposte accurate e
        un'interfaccia user-friendly.
        \item Memoria a lungo termine/Salvataggio chat recenti.
        \item Velocità di risposta e disponibilità 24/7.
    \end{itemize}

    \subsection{Ambito del Progetto}

        \subsubsection{Inclusioni}
        \begin{itemize}
            \item Database relazionale per la gestione dei dati sui prodotti.
            \item Integrazione di un modello \href{https://code7crusaders.github.io/docs/PB/documentazione_interna/glossario.html#llm-large-language-model}{LLM\textsuperscript{G}} tramite API.
            \item Interfaccia utente mobile per l'interazione con l'IA.
            \item Funzionalità di configurazione backend per template di domande e risposte.
        \end{itemize}
        \subsubsection{Esclusioni}
        \begin{itemize}
            \item Creazione di un nuovo modello \href{https://code7crusaders.github.io/docs/PB/documentazione_interna/glossario.html#llm-large-language-model}{LLM\textsuperscript{G}}.
            \item Supporto a lingue non previste dal modello \href{https://code7crusaders.github.io/docs/PB/documentazione_interna/glossario.html#llm-large-language-model}{LLM\textsuperscript{G}} scelto.
        \end{itemize}

    \subsection{Funzionamento}
    \begin{itemize}
        \item Da un’interfaccia utente, viene catturata una domanda da parte dell’utente.
        \item La domanda viene inoltrata al sistema attraverso delle \href{https://code7crusaders.github.io/docs/PB/documentazione_interna/glossario.html#api-rest-representational-state-transfer}{API REST\textsuperscript{G}} risiedenti in un Web Server.
        \item La query ricevuta viene gestita dall’\href{https://code7crusaders.github.io/docs/PB/documentazione_interna/glossario.html#embedding}{Embedding\textsuperscript{G}} Model che trasforma la domanda in rappresentazione vettoriale.
        \item La rappresentazione vettoriale viene utilizzata per effettuare una ricerca all’interno del database vettoriale da dove vengono reperiti i vettori più simili
        \item Sia la domanda sia i risultati della ricerca nel database vettoriale, vengono inviati al sistema \href{https://code7crusaders.github.io/docs/PB/documentazione_interna/glossario.html#llm-large-language-model}{LLM\textsuperscript{G}} che costruirà la risposta utilizzando il contesto fornito.
        \item Attraverso \href{https://code7crusaders.github.io/docs/PB/documentazione_interna/glossario.html#api-rest-representational-state-transfer}{API REST\textsuperscript{G}}, il sistema inoltra la riposta al dispositivo dell’utente.
    \end{itemize}

    \subsection{Tecnologie e Strumenti Consigliati}
    L’azienda proponente è disponibile a fornire i dati di un caso di studio da utilizzare
    per lo sviluppo del progetto. I dati potranno essere dati in ingresso al sistema così da eseguire la fase di
    training e poi interagire con il sistema per valutarne le prestazioni sfruttando un caso reale.
    Di seguito vengono suggerite alcune tecnologie utilizzabili per il sistema esposto:
    \begin{itemize}
        \item \textbf{Database}: MySQL o \href{https://code7crusaders.github.io/docs/PB/documentazione_interna/glossario.html#postgresql}{PostgreSQL\textsuperscript{G}}.
        \item \href{https://code7crusaders.github.io/docs/PB/documentazione_interna/glossario.html#llm-large-language-model}{\textbf{LLM}\textsuperscript{G}}: BLOOM o Italia by iGenius, in base alle prestazioni richieste.
        \item \href{https://code7crusaders.github.io/docs/PB/documentazione_interna/glossario.html#backend}{\textbf{Backend}\textsuperscript{G}}: Node.js con Express.js o .NET.
        \item \href{https://code7crusaders.github.io/docs/PB/documentazione_interna/glossario.html#frontend}{\textbf{Frontend}\textsuperscript{G}}: .NET MAUI per applicazioni mobile multipiattaforma.
        \item \href{https://code7crusaders.github.io/docs/PB/documentazione_interna/glossario.html#api-rest-representational-state-transfer}{\textbf{API REST}\textsuperscript{G}}: Per la comunicazione tra \href{https://code7crusaders.github.io/docs/PB/documentazione_interna/glossario.html#llm-large-language-model}{LLM\textsuperscript{G}} e interfaccia utente.
        \item \textbf{Controllo Versione}: Git (GitHub per repository pubblico).
    \end{itemize}

%%%%%%%%%%%%%%%%%%%%%%%%%%%%%%%%%%%%%%%%%%%%%%%%%%%%%%%%%%%%%%%%%%%%%%%%%%%%%%%%%%%%%%



% Pianificazione %%%%%%%%%%%%%%%%%%%%%%%%%%%%%%%%%%%%%%%%%%%%%%%%%%%%%%%%%%%%%%%%%%%%%%%%
\newpage
\section{Pianificazione}
    \subsection{Struttura del Team}
        \subsubsection{Ruoli}
        I ruoli in seguito descritti sono equamente divisi tra i vari componenti del Team. Ogni ruolo possiede diversi incarichi e obbiettivi:
        \begin{itemize}
            \item \textbf{Responsabile}: coordina il gruppo di lavoro, controlla le attività e gestisce le risorse. Si occupa di garantire che il progetto venga portato a termine nei tempi stabiliti e con le risorse disponibili.
            \item \textbf{Amministratore}: si occupa della gestione delle risorse e delle infrastrutture, incluso il setup degli strumenti di supporto alla produzione del software. Garantisce inoltre l’uso corretto delle procedure per assicurare efficienza e produttività.
            \item \textbf{Analista}: gioca un ruolo fondamentale nella fase iniziale del progetto. È responsabile della definizione dei requisiti e dell’analisi delle funzionalità del software, delineando i casi d'uso. Essendo necessario principalmente all'inizio del progetto, il numero di ore assegnato al ruolo è relativamente ridotto.
            \item \textbf{Progettista}: definisce l'architettura del software, descrivendo le componenti e le loro interazioni sulla base dei requisiti stabiliti dall'Analista. Questo ruolo ha un numero di ore significativamente elevato perché è essenziale per garantire una struttura solida, soprattutto considerando l’implementazione di modelli \href{https://code7crusaders.github.io/docs/PB/documentazione_interna/glossario.html#llm-large-language-model}{\emph{LLM}\textsuperscript{G}}, che richiedono un'architettura ben progettata e adattata a tali tecnologie.
            \item \textbf{Programmatore}: si occupa di scrivere il codice del software seguendo le specifiche del progettista. Il numero di ore assegnato è alto, dato che rappresenta il cuore della fase di sviluppo. Tuttavia, il ruolo ha leggermente meno ore rispetto al Verificatore, poiché abbiamo scelto di adottare una metodologia incentrata sui test, che richiede un’accurata verifica del software.
            \item \textbf{Verificatore}: verifica che il software e la documentazione siano conformi alle norme e alle specifiche. Questo ruolo richiede un numero di ore superiore alla media, data la necessità di test approfonditi e continui, in particolare per un progetto basato su \href{https://code7crusaders.github.io/docs/PB/documentazione_interna/glossario.html#llm-large-language-model}{\emph{LLM}\textsuperscript{G}}, dove ogni componente deve essere rigorosamente validato per garantire la precisione e l’affidabilità del sistema.
        \end{itemize}
        \subsubsection{Stakeholder}
        \begin{itemize}
            \item \textbf{Cliente}: Ergon Informatica Srl.
            \item \textbf{Referente interno}: Gianluca Carlesso.
        \end{itemize}

    \subsection{Budget e Risorse}
        \subsubsection{TODO: }
        \begin{itemize}
            \item \textbf{Allocazione hardware}: Server dedicati per il database e le API, oltre a risorse cloud per l'hosting del modello \href{https://code7crusaders.github.io/docs/PB/documentazione_interna/glossario.html#llm-large-language-model}{LLM\textsuperscript{G}}.
            \item \textbf{Licenze software e costi}: Costi delle licenze software necessarie e del modello \href{https://code7crusaders.github.io/docs/PB/documentazione_interna/glossario.html#llm-large-language-model}{LLM\textsuperscript{G}}, se applicabili.
        \end{itemize}
        \subsubsection{Distribuzione ore/ruolo}
        Di seguito, si riporta il costo orario in base al ruolo assunto:
        \begin{table}[!h]
            \begin{center}
                \begin{tabular}{ |c|c|c|c| }
                    \hline
                    \textbf{Ruolo}          & \textbf{Costo orario} (\euro) &  \textbf{per ruolo}   & \textbf{Ore per membro} \\
                    \hline
                    Responsabile   & 30           &     54       &       8        \\
                    Amministratore & 20           &     64       &       9        \\
                    Analista       & 25           &     65       &       9       \\
                    Progettista    & 25           &     105      &       15       \\
                    Programmatore  & 15           &     184      &       26       \\
                    Verificatore   & 15           &     193      &       28       \\
                    \hline
                    \textbf{Totale}         &    12805    &     665       &       95       \\
                    \hline
                \end{tabular}
                \caption{Costo orario e totale}
            \end{center}
        \end{table}

        \subsubsection{Distribuzione ore/membro}
        Tutti i componenti del Team Code7Crusaders si impegnano a dedicare un totale di \textbf{95 ore} di lavoro effettivo partizionate
        settimanalmente in base al ruolo di riferimento, per lo svolgimento del capitolato \textbf{C7} di \textbf{Ergon Informatica}. Inoltre,
        ciascun membro garantisce la conclusione del progetto entro la data prevista e preventivata nel paragrafo 5 di questo documento.
        \\
        Ripartizione delle ore per membro del team:
        \begin{table}[!h]
            \begin{center}
                \begin{tabular}{ |c|c|c|c|c|c|c|c| }
                    \hline
                    \textbf{Membro}    & \textbf{Re} & \textbf{Am} & \textbf{An} & \textbf{Pj} & \textbf{Pg} & \textbf{Ve} & \textbf{Totale} \\
                    \hline
                    Enrico Cotti Cottini     & 8           & 9           & 9          & 15          & 26          & 28          & 95              \\
                    Gabriele Di Pietro       & 8           & 9           & 9          & 15          & 26          & 28          & 95              \\
                    Tommaso Diviesti         & 8           & 9           & 9          & 15          & 26          & 28          & 95              \\
                    Francesco Lapenna        & 8           & 9           & 9          & 15          & 26          & 28          & 95              \\
                    Matthew Pan              & 8           & 9           & 9          & 15          & 26          & 28          & 95              \\
                    Eddy Pinarello           & 8           & 9           & 9          & 15          & 26          & 28          & 95              \\
                    Filippo Rizzolo          & 8           & 9           & 9          & 15          & 26          & 28          & 95              \\
                    \hline
                \end{tabular}
                \caption{Impegni orari a persona}
            \end{center}
        \end{table}
        \\
        \textsc{Legenda:} \\
            \textbf{Re} = Responsabile \\
            \textbf{Am} = Amministratore \\
            \textbf{An} = Analista \\
            \textbf{Pj} = Progettista \\
            \textbf{Pg} = Programmatore \\
            \textbf{Ve} = Verificatore \\
        
            \subsection{Piano di Comunicazione}
            \begin{itemize}
                \item Riunioni settimanali con il team.
                \item Report di avanzamento per il referente aziendale ogni 2 settimane.
                \item Feedback continuo attraverso test intermedi.
            \end{itemize}
        
            \subsection{Pianificazione delle Attività}
                Il gruppo Code7Crusaders
                si impegna a consegnare il progetto entro il 14/03/2025. La pianificazione prevede
                19 settimane di lavoro, suddivise come segue:
                \begin{itemize}
                    \item \textbf{Sviluppo POC (\textit{Proof of Concept}): 6 settimane}
                    \item \textbf{Sviluppo MVP (\textit{Minimum Viable Product}): 13 settimane}
                \end{itemize}
        
                \subsubsection{Fasi principali}  % TODO: rivedere le fasi
                \begin{enumerate}
                    \item \textbf{Analisi dei requisiti} (2 settimane):
                    \begin{itemize}
                        \item Revisione del capitolato.
                        \item Identificazione delle tecnologie e dei modelli \href{https://code7crusaders.github.io/docs/PB/documentazione_interna/glossario.html#llm-large-language-model}{LLM\textsuperscript{G}} adatti.
                    \end{itemize}
                    \item \textbf{Progettazione} (2 settimane):
                    \begin{itemize}
                        \item Progettazione architetturale.
                        \item Definizione dello schema del database.
                    \end{itemize}
                    \item \textbf {Sviluppo Backend} (4 settimane):
                    \begin{itemize}
                        \item Configurazione del database.
                        \item Implementazione delle API REST.
                    \end{itemize}
                    \item \textbf{Integrazione LLM} (3 settimane):
                    \begin{itemize}
                        \item Pre-processing dei dati e integrazione del modello \href{https://code7crusaders.github.io/docs/PB/documentazione_interna/glossario.html#llm-large-language-model}{LLM\textsuperscript{G}}.
                    \end{itemize}
                    \item \textbf{Sviluppo Frontend} (4 settimane):
                    \begin{itemize}
                        \item Creazione dell'interfaccia utente mobile.
                    \end{itemize}
                    \item \textbf{Test e validazione} (2 settimane):
                    \begin{itemize}
                        \item Test funzionali e di usabilità.
                    \end{itemize}
                    \item \textbf{Rilascio e documentazione} (2 settimane).
                \end{enumerate}
        
                \subsubsection{Cronoprogramma}
                % Diagramma di Gantt %%%%%%%%%%%%%%%%%%%%%%%%%%%%%%%%%%%%%%%%%%%%%%%%%%%%%%%%%%%%
                \resizebox{\textwidth}{!}{ % Scales the chart to fit the text width
                    \begin{ganttchart}[
                        hgrid, vgrid, % Adds grid lines
                        x unit=1.1cm, % Adjust the width of each time slot
                        title/.append style={draw=black, thick, fill=red!20, line width=1pt, draw opacity=1},
                        bar/.append style={fill=red!50}, % Custom bar styling
                        milestone/.append style={fill=red!70}, % \href{https://code7crusaders.github.io/docs/PB/documentazione_interna/glossario.html#milestone}{Milestone\textsuperscript{G}} style
                        ]{1}{19} % Timeline from week 1 to week 19
                        \gantttitle{Project Plan (Weeks)}{19} \\ % Title spanning 19 weeks
                        \gantttitle{Novembre 2024}{4}
                        \gantttitle{Dicembre 2024}{5}
                        \gantttitle{Gennaio 2025}{4}
                        \gantttitle{Febbraio 2025}{4}
                        \gantttitle{Marzo 2025}{2}
                        \\
                        \gantttitle{4 nov}{1}  % TODO: decidere cosa tenere tra i titoli
                        \gantttitle{11 nov}{1}
                        \gantttitle{18 nov}{1}
                        \gantttitle{25 nov}{1}
                        \gantttitle{2 dic}{1}
                        \gantttitle{9 dic}{1}
                        \gantttitle{16 dic}{1}
                        \gantttitle{23 dic}{1}
                        \gantttitle{30 dic}{1}
                        \gantttitle{6 gen}{1}
                        \gantttitle{13 gen}{1}
                        \gantttitle{20 gen}{1}
                        \gantttitle{27 gen}{1}
                        \gantttitle{3 feb}{1}
                        \gantttitle{10 feb}{1}
                        \gantttitle{17 feb}{1}
                        \gantttitle{24 feb}{1}
                        \gantttitle{3 mar}{1}
                        \gantttitle{10 mar}{1} \\
                        \gantttitlelist{1,...,19}{1} \\ % Weekly columns
        
                        % Define groups and tasks
                        \ganttgroup{Sviluppo PoC}{1}{6} \\
                        \ganttmilestone{RTB}{6} \\
                        \ganttgroup{Sviluppo MVP}{7}{19} \\
                        \ganttmilestone{PB}{19} \\
        
                        \ganttbar{Analisi dei requisiti}{1}{2} \\
                        \ganttbar{Progettazione}{3}{4} \\
                        \ganttbar{Sviluppo backend}{5}{8} \\
                        \ganttbar{Integrazione LLM}{9}{11} \\
                        \ganttbar{Sviluppo frontend}{12}{15} \\
                        \ganttbar{Test e validazione}{16}{17} \\
                        \ganttbar{Rilascio e documentazione}{18}{19} \\
        
                        \ganttlink{elem4}{elem5}  % TODO: decidere se tenere o togliere i link
                        \ganttlink{elem5}{elem6}
                        \ganttlink{elem6}{elem7}
                        \ganttlink{elem7}{elem8}
                        \ganttlink{elem8}{elem9}
                        \ganttlink{elem9}{elem10}
        
                    \end{ganttchart}
                }
        
    \section{Analisi dei Rischi}
    \subsection{Introduzione}
    In questa sezione vengono analizzati i principali rischi che potrebbero manifestarsi durante lo svolgimento del progetto, con una valutazione della loro gravità e probabilità di occorrenza. Ad ogni rischio è associato un indice numerico che ne definisce l'intensità e la probabilità, consentendo di determinare la criticità di ciascun rischio.

    \subsection{Definizione degli Indici}
        Per una corretta valutazione dei rischi, sono stati definiti i seguenti fattori chiave:
        \begin{itemize}
            \item \textbf{Indice}: un valore numerico progressivo che identifica in modo univoco ogni rischio;
            \item \textbf{Tipo}: rappresenta la categoria di impatto che il rischio potrebbe avere sul progetto, ed è classificato in Basso, Medio o Alto;
            \item \textbf{Gravità}: descrive le conseguenze o l'impatto che un rischio avrebbe se si verificasse, con l'eventuale necessità di modificare la pianificazione o l'approccio al progetto;
            \item \textbf{Probabilità}: indica la possibilità che il rischio si verifichi durante il ciclo di vita del progetto.
        \end{itemize}
        Gli indici di \textbf{Gravità} e di \textbf{Probabilità} sono definiti con i seguenti valori:
        \begin{table}[h!]
            \centering
            \begin{tabular}{|c|c|p{6cm}|p{6cm}|}
                \hline
                \textbf{Indice} & \textbf{Tipo} & \textbf{Gravità} & \textbf{Probabilità} \\
                \hline
                \textbf{1} & Basso & L'impatto sul progetto è minimo o trascurabile, ad esempio un lieve rallentamento senza effetti sui tempi di consegna & La possibilità che si verifichi è bassa, ma esistono fattori che potrebbero contribuire alla sua realizzazione \\
                \hline
                \textbf{2} & Medio & L'evento richiede l'uso di risorse aggiuntive o un parziale adattamento del piano, con impatti gestibili ma che richiedono uno sforzo supplementare & Esiste una probabilità media che il rischio si realizzi \\
                \hline
                \textbf{3} & Alto & Il rischio comporta ritardi significativi, aumento dei costi o degrado della qualità, impattando negativamente sull'esperienza utente e sul raggiungimento degli obiettivi del progetto & È altamente probabile che il rischio si verifichi, con segni evidenti che ne indicano la possibilità \\
                \hline
            \end{tabular}
            \caption{Definizione degli Indici di Gravità e Probabilità}
            \label{tab:definizione_indici}
        \end{table}
            
    \subsection{Suddivisione dei rischi}
    I rischi del progetto sono suddivisi in tre categorie principali, in base alla natura delle problematiche che potrebbero verificarsi:
    \begin{itemize}
        \item Rischi di natura tecnologica;
        \item Rischi legati alla comunicazione;
        \item Rischi relativi all'organizzazione e gestione del progetto.
        \item Rischi individuali
    \end{itemize}
    Ogni rischio è identificato tramite un codice univoco, strutturato come segue:
    \[
    \textbf{R[Categoria][Indice] - [Nome]}
    \]
    Dove:
    \begin{itemize}
        \item \textbf{Categoria}: rappresenta la tipologia del rischio, che può essere:
        \begin{itemize}
            \item \textbf{T}: per i rischi tecnologici;
            \item \textbf{C}: per i rischi comunicativi;
            \item \textbf{P}: per i rischi di pianificazione;
            \item \textbf{I}: per i rischi individuali;
        \end{itemize}
        \item \textbf{Indice}: è un identificativo progressivo univoco all'interno della categoria.
        \item \textbf{Nome}: è una breve descrizione del rischio.
    \end{itemize}

    \subsection{Rischi Tecnologici}

    \subsubsection{RT1 - Difficoltà	 nell’uso	di nuove tecnologie}
    \begin{table}[H]
        \renewcommand{\arraystretch}{1.5}
        \centering
        \begin{tabular}{|p{3cm}|p{11cm}|}
            \hline
            \textbf{Codice} & RT1 \\
            \hline
            \textbf{Descrizione} & Rischio legato all’inesperienza o alla poca conoscenza di un componente necessario da integrare all’interno del progetto. Questo potrebbe causare ritardi sia nella fase di progettazione che in quella di sviluppo. \\
            \hline
            \textbf{Gravità} & 2 \\
            \hline
            \textbf{Probabilità} & 3 \\
            \hline
            \textbf{Rilevamento} & Comunicazione delle proprie conoscenze e difficoltà. \\
            \hline
            \textbf{Contromisure} & Il team si impegnerà a studiare in modo approfondito le tecnologie richieste dal capitolato, in particolare le tecnologie di \href{https://code7crusaders.github.io/docs/PB/documentazione_interna/glossario.html#llm-large-language-model}{LLM\textsuperscript{G}}. Verranno organizzati corsi di formazione interna per permettere ai membri di acquisire competenze comuni e supportarsi vicendevolmente. \\
            \hline
        \end{tabular}
        \caption{Tabella RT1}
    \end{table}

    \subsubsection{RT2 - Sviluppo di codice poco comprensibile}
    \begin{table}[H]
        \renewcommand{\arraystretch}{1.5}
        \centering
        \begin{tabular}{|p{3cm}|p{11cm}|}
            \hline
            \textbf{Codice} & RT2 \\
            \hline
            \textbf{Descrizione} & Durante il processo di sviluppo, un membro del Team potrebbe scrivere una o più porzioni di codice difficile da comprendere per tutti i membri. In questo caso, si perde del tempo che potrebbe essere utilizzato per svolgere altri compiti. \\
            \hline
            \textbf{Gravità} & 2 \\
            \hline
            \textbf{Probabilità} & 2 \\
            \hline
            \textbf{Rilevamento} & Revisione del codice scritto da ogni programmatore da parte di un verificatore. \\
            \hline
            \textbf{Contromisure} & Effettuare revisioni regolari del codice. Necessaria, inoltre, una comunicazione tra gli sviluppatori in modo che il processo di sviluppo sia facilitato e la struttura sia comprensibile a tutti. \\
            \hline
        \end{tabular}
        \caption{Tabella RT2}
    \end{table}

    
    \subsection{Rischi comunicativi}
            
    \subsubsection{RC1 - Disaccordi all’interno del Team}
    \begin{table}[H]
        \renewcommand{\arraystretch}{1.5}
        \centering
        \begin{tabular}{|p{3cm}|p{11cm}|}
            \hline
            \textbf{Codice} & RC1 \\
            \hline
            \textbf{Descrizione} & I componenti del Team possono avere ideologie e opinioni diverse riguardo elementi particolari del progetto. \\
            \hline
            \textbf{Gravità} & 2 \\
            \hline
            \textbf{Probabilità} & 1 \\
            \hline
            \textbf{Rilevamento} & Possono essere identificate attraverso le opinioni espresse o osservando le dinamiche del gruppo. \\
            \hline
            \textbf{Contromisure} & In caso di disaccordo, si procederà con una discussione interna dove si elencheranno i punti a favore e i contro in modo da arrivare ad una soluzione comune. \\
            \hline
        \end{tabular}
        \caption{Tabella RC1}
    \end{table}

    \subsubsection{RC2 - Problemi di comunicazione all’interno del gruppo}
    \begin{table}[H]
        \renewcommand{\arraystretch}{1.5}
        \centering
        \begin{tabular}{|p{3cm}|p{11cm}|}
            \hline
            \textbf{Codice} & RC2 \\
            \hline
            \textbf{Descrizione} & Situazioni in cui non esistono linee guida chiare per la trasmissione efficace delle informazioni all’interno del team. Questa mancanza di direzione può portare a fraintendimenti, confusione o ritardi. \\
            \hline
            \textbf{Gravità} & 2 \\
            \hline
            \textbf{Probabilità} & 1 \\
            \hline
            \textbf{Rilevamento} & Identificati tramite sondaggi, feedback e comportamenti da parte dei membri del gruppo durante le riunioni o comunicazioni via messaggio. \\
            \hline
            \textbf{Contromisure} & Oltre ad aver stabilito i canali specifici di comunicazione per i vari confronti, sono fissate delle riunioni settimanali dove vengono discusse tutte le questioni all’ordine del giorno. Inoltre, in caso di necessità, si possono pianificare dei meeting straordinari per parlare di questioni imminenti ed importanti. \\
            \hline
        \end{tabular}
        \caption{Tabella RC2}
    \end{table}

    \subsubsection{RC3 - Mancato confronto con l’azienda Proponente}
    \begin{table}[H]
        \renewcommand{\arraystretch}{1.5}
        \centering
        \begin{tabular}{|p{3cm}|p{11cm}|}
            \hline
            \textbf{Codice} & RC3 \\
            \hline
            \textbf{Descrizione} & Questo rischio identifica una mancanza di comunicazione con l’azienda Proponente che impedisce l’acquisizione di chiarimenti durante lo svolgimento del progetto. \\
            \hline
            \textbf{Gravità} & 3 \\
            \hline
            \textbf{Probabilità} & 1 \\
            \hline
            \textbf{Rilevamento} & Implicito \\
            \hline
            \textbf{Contromisure} & Cercare di mantenere l’azienda Proponente informata sugli sviluppi in modo regolare e chiedere feedback quando necessario. \\
            \hline
        \end{tabular}
        \caption{Tabella RC3}
    \end{table}

    \break
    \subsection{Rischi organizzattivi}
    
    \subsubsection{RO1 - Incomprensioni nella pianificazione delle attività}
    \begin{table}[H]
        \renewcommand{\arraystretch}{1.5}
        \centering
        \begin{tabular}{|p{3cm}|p{11cm}|}
            \hline
            \textbf{Codice} & RO1 \\
            \hline
            \textbf{Descrizione} & Si verifica quando il Team manca di esperienza o competenza nella gestione dei processi di pianificazione, portando a possibili conseguenze come l’allocazione inefficace delle risorse, la scarsa definizione dei requisiti o valutazioni inaccurate dei tempi. \\
            \hline
            \textbf{Gravità} & 3 \\
            \hline
            \textbf{Probabilità} & 2 \\
            \hline
            \textbf{Rilevamento} & Analisi di progresso/tempo e aderenza alle scadenze. \\
            \hline
            \textbf{Contromisure} & Nel caso in cui uno o più componenti del gruppo abbiano difficoltà nello svolgimento di determinati compiti previsti, anche in caso di scadenze imminenti, verrà effettuata una revisione delle attività per dare priorità a quelle con scadenze in direzione d’arrivo. \\
            \hline
        \end{tabular}
        \caption{Tabella RO1}
    \end{table}

    \subsubsection{RO2 - Deviazione dai tempi e costi previsti}
    \begin{table}[H]
        \renewcommand{\arraystretch}{1.5}
        \centering
        \begin{tabular}{|p{3cm}|p{11cm}|}
            \hline
            \textbf{Codice} & RO2 \\
            \hline
            \textbf{Descrizione} & Questo rischio si riferisce alla possibilità di variazioni significative rispetto alle tempistiche e ai costi pianificati inizialmente per lo sviluppo del progetto. \\
            \hline
            \textbf{Gravità} & 3 \\
            \hline
            \textbf{Probabilità} & 1 \\
            \hline
            \textbf{Rilevamento} & Monitoraggio di tempi e costi attraverso il consuntivo periodico (sprint) presente nel piano di progetto. \\
            \hline
            \textbf{Contromisure} & Il team monitorerà costantemente l’avanzamento del progetto attraverso riunioni settimanali di aggiornamento e controllo del budget. Eventuali scostamenti saranno analizzati per prendere tempestivamente le misure correttive. \\
            \hline
        \end{tabular}
        \caption{Tabella RO2}
    \end{table}

    \break
    \subsection{Rischi individuali}

    \subsubsection{RI1 - Scarsa collaborazione da parte di uno o più membri del team}
    \begin{table}[H]
        \renewcommand{\arraystretch}{1.5}
        \centering
        \begin{tabular}{|p{3cm}|p{11cm}|}
            \hline
            \textbf{Codice} & RI1 \\
            \hline
            \textbf{Descrizione} & La possibilità che uno o più membri del gruppo non collaborino attivamente allo sviluppo del progetto. \\
            \hline
            \textbf{Gravità} & 3 \\
            \hline
            \textbf{Probabilità} & 1 \\
            \hline
            \textbf{Rilevamento} & Implicito. \\
            \hline
            \textbf{Contromisure} & Saranno programmati incontri regolari di coordinamento per favorire la comunicazione e la collaborazione all’interno del team. Ogni membro avrà ruoli e responsabilità ben definiti, e verranno incentivati il dialogo aperto e la trasparenza. \\
            \hline
        \end{tabular}
        \caption{Tabella RI1}
    \end{table}

    \subsubsection{RI2 - Impegni personali e universitari }
    \begin{table}[H]
        \renewcommand{\arraystretch}{1.5}
        \centering
        \begin{tabular}{|p{3cm}|p{11cm}|}
            \hline
            \textbf{Codice} & RI2 \\
            \hline
            \textbf{Descrizione} & Questo rischio riguarda la potenziale difficoltà dei membri del team, nel dedicare il tempo richiesto al progetto a causa degli impegni universitari e/o altri impegni personali. \\
            \hline
            \textbf{Gravità} & 1 \\
            \hline
            \textbf{Probabilità} & 2 \\
            \hline
            \textbf{Rilevamento} & Comunicazione della propria disponibilità. \\
            \hline
            \textbf{Contromisure} & Ogni membro del Team dovrà segnalare in anticipo eventuali impegni personali o universitari. Il team utilizzerà un calendario condiviso per pianificare le attività in modo da evitare sovrapposizioni e garantire la disponibilità necessaria. \\
            \hline
        \end{tabular}
        \caption{Tabella RI2}
    \end{table}

    \break
    \subsection{Tabella riassuntiva}
    \begin{table}[H]
        \renewcommand{\arraystretch}{1.5}
        \centering
        \begin{tabular}{|p{1.5cm}|p{10cm}|p{2cm}|p{2cm}|}
            \hline
            \textbf{Codice} & \textbf{Rischio} & \textbf{Gravità} & \textbf{Probabilità} \\
            \hline
            RT1 & Difficoltà nell’uso di nuove tecnologie & Media & Alta \\
            \hline
            RT2 & Sviluppo di codice poco comprensibile & Media & Media \\
            \hline
            RC1 & Disaccordi all’interno del Team & Media & Bassa \\
            \hline
            RC2 & Problemi di comunicazione all’interno del gruppo & Media & Bassa \\
            \hline
            RC3 & Mancato confronto con l’azienda Proponente & Alta & Bassa \\
            \hline
            RO1 & Incomprensioni nella pianificazione delle attività & Alta & Media \\
            \hline
            RO2 & Deviazione dai tempi e costi previsti & Alta & Bassa \\
            \hline
            RI1 & Scarsa collaborazione da parte di uno o più membri del team & Alta & Bassa \\
            \hline
            RI2 & Impegni personali e universitari  & Bassa & Media \\
            \hline
        \end{tabular}
        \caption{Tabella riassuntiva rischi individuati}
    \end{table}


        %%%%%%%%%%%%%%%%%%%%%%%%%%%%%%%%%%%%%%%%%%%%%%%%%%%%%%%%%%%%%%%%%%%%%%%%%%%%%%%%%



%%%%%%%%%%%%%%%%%%%%%%%%%%%%%%%%%%%%%%%%%%%%%%%%%%%%%%%%%%%%%%%%%%%%%%%%%%%%%%%%%%%%%%
\newpage
\section{Gestione del modello di sviluppo}
    \subsection{Modello adottato}
    Dopo aver effettuato una analisi delle esigenze del progetto e una valutazione
    di gruppo, il Team ha deciso di adottare il modello \textbf{Agile}. È stato scelto il 
    modello Agile perché ci consente di lavorare in modo più flessibile e reattivo 
    alle esigenze in continua evoluzione del progetto e dei clienti. Grazie alla sua 
    iterazione continua e al focus sulla collaborazione, possiamo migliorare rapidamente 
    il prodotto, ottenere feedback costante e adattarci prontamente ai cambiamenti. 
    I periodi di sviluppo sono divisi in \textbf{sprint settimanali} che sono accompagnati da riunioni
    periodiche tra i vari componenti del Team e, quando necessario, anche con l'Azienda proponente.
    Questo serve anche per promuove una comunicazione trasparente all'interno del team, 
    migliorando l'efficienza e la qualità del lavoro.
    \subsection{Vantaggi del modello}
    \begin{itemize}
        \item \textbf{Flessibilità e adattabilità}: permette di adattarsi rapidamente ai cambiamenti. Poiché i progetti sono suddivisi in iterazioni brevi (sprint), è possibile modificare il piano di lavoro in corso d'opera in base ai feedback e alle nuove informazioni.
        \item \textbf{Rilasci frequenti e valore continuo}: con la consegna di piccole porzioni di prodotto a intervalli regolari si ottengono maggiori feedback e gli errori vengono corretti più velocemente dato che i risultati sono monitorati e valutati frequentemente.
        \item \textbf{Collaborazione tra il team}: la metodologia Agile promuove il lavoro collaborativo tra membri del team con competenze diverse. Questo aiuta a ottenere soluzioni più complete e ben integrate, che rispondano meglio alle necessità del prodotto.
        \item \textbf{Migliore gestione dei rischi}: grazie alla continua revisione dei progressi e al rilascio di versioni parziali del prodotto, è possibile identificare e affrontare i problemi in modo tempestivo.
        \item \textbf{Maggiore qualità del prodotto}: con test continui e revisioni regolari del codice e del prodotto, si migliora la qualità del software o del prodotto finale. Le iterazioni frequenti consentono di correggere rapidamente eventuali errori e migliorare il prodotto in modo incrementale.
        \item \textbf{Migliore gestione delle risorse e tempi di consegna}: l'approccio iterativo permette di pianificare e gestire le risorse in modo più efficiente. Poiché il lavoro è suddiviso in piccoli blocchi, è più facile stimare tempi e risorse necessari, evitando sovraccarichi di lavoro e garantendo una consegna puntuale.
        \item \textbf{Incremento della produttivita}: le funzionalità più importanti o urgenti vengono sviluppate prima, concentrando risorse e sforzi sulle attività che generano maggior valore. Ciò ottimizza l'uso del tempo e delle risorse, aumentando la produttività complessiva.
    \end{itemize}
    \subsection{Periodi}
    Ogni periodo è caratterizzato dai seguenti elementi:
    \begin{itemize}
        \item Data di inizio, data di fine prevista, data di fine effettiva ed eventuali giorni di ritardo.
        \item Pianificazione delle attivita da svolgere con possibili rischi.
        \item Stima temporale al completamento di tutte le attività previste.
        \item Paragone temporale tra lavoro svolto e preventivato.
        \item Prospetto economico relativo ai ruoli svolti.
        \item Valutazione impatto e mitigazione dei rischi effettivamente occorsi.
        \item Retrospettiva per l'automiglioramento.
    \end{itemize}

    % TODO Sistemare gli sprint con le caratteristiche elencate qui sopra

        \subsection{Requirements and Technology Baseline (RTB)}
            \subsubsection{Sprint 1}
                % Sprint Planning
                \paragraph{Pianificazione}
                \begin{itemize}
                    \item inizio: 4 Novembre 2024
                    \item fine: 10 Novembre 2024
                \end{itemize}
                
                \paragraph{Obiettivi dello Sprint}
                    \begin{itemize}
                        \item Redigere il documento \href{https://code7crusaders.github.io/docs/PB/documentazione_interna/glossario.html#norme-di-progetto}{\emph{Norme di Progetto}\textsuperscript{G}}.
                        \item Approfondire l'uso delle tecnologie richieste nel capitolato.
                        \item Effettuare una chiamata con l'azienda per definire i requisiti utente e software.
                        \item Studiare e comprendere i documenti legati alla documentazione.
                        \item Migliorare il sito web aggiungendo un glossario interattivo.
                    \end{itemize}

                \paragraph{Task Prioritari}
                    \begin{itemize}
                        \item Automazione del processo di compilazione e deploy dei file LaTeX.
                        \item Organizzazione degli incontri settimanali per assegnare i compiti.
                        \item Creazione di workflow per la gestione efficiente del progetto.
                    \end{itemize}

                \paragraph{Rischi attesi}\mbox{}\\
                I componenti del Team ritengono che per questo Sprint siano possibili i seguenti rischi:
                \begin{itemize}
                    \item RT1 - Difficoltà nell’uso di nuove tecnologie
                    \item RO1 - Incomprensioni nella pianificazione delle attività
                    \item RI2 - Impegni personali e universitari
                \end{itemize}

                % TODO sposta i ruoli in base allo sprint
                \paragraph{Preventivo}\mbox{}\\
                \begin{table}[H]
                    \centering
                    \begin{tabular}{|c|c|c|c|c|c|c|c|}
                    \hline
                                  & Re  & Am  & An  & Pj  & Pg  & Ve  & Totale per persona \\ \hline
                    Cotti Cottini & 3   & -   & -   & -   & -   & -   & 3                  \\ \hline
                    Di Pietro     & -   & -   & -   & 5,5 & -   & -   & 5,5                \\ \hline
                    Diviesti      & -   & 3,5 & -   & -   & -   & -   & 3,5                \\ \hline
                    Lapenna       & -   & -   & 3,5 & -   & -   & -   & 3,5                \\ \hline
                    Pan           & -   & -   & -   & -   & -  & 4   & 4                 \\ \hline
                    Pinarello     & -   & -   & -   & -   & -   & 4   & 4                  \\ \hline
                    Rizzolo       & -   & -   & -   & -   & -   & 4   & 4                  \\ \hline
                    Totale        & 3   & 3,5 & 3,5 & 5,5 & -  & 12  &                    \\ \hline
                    \end{tabular}
                    \caption{Preventivo orario per ruolo nello Sprint 1}
                \end{table}


                \begin{figure}[H]
                    \centering
                    \begin{tikzpicture}
                        \pie[polar]{11/Responsabile, 13/Amministratore, 13/Analista, 20/Progettista, 43/Verificatore}
                    \end{tikzpicture}
                    \caption{Ripartizione in percentuale dei ruoli nello Sprint 1}
                \end{figure}


                % TODO copiala dal preventivo (per ora li facciamo uguali)
                \paragraph{Consuntivo}\mbox{}\\
                \begin{table}[H]
                    \centering
                    \begin{tabular}{|c|c|c|c|c|c|c|c|}
                        \hline
                                      & Re  & Am  & An  & Pj  & Pg  & Ve  & Totale per persona \\ \hline
                        Cotti Cottini & 3   & -   & -   & -   & -   & -   & 3                  \\ \hline
                        Di Pietro     & -   & -   & -   & 5,5 & -   & -   & 5,5                \\ \hline
                        Diviesti      & -   & 3,5 & -   & -   & -   & -   & 3,5                \\ \hline
                        Lapenna       & -   & -   & 3,5 & -   & -   & -   & 3,5                \\ \hline
                        Pan           & -   & -   & -   & -   & -  & 4   & 4                 \\ \hline
                        Pinarello     & -   & -   & -   & -   & -   & 5 \textcolor{red}{(+1)}  & 5 \textcolor{red}{(+1)}                 \\ \hline
                        Rizzolo       & -   & -   & -   & -   & -   & 5 \textcolor{red}{(+1)}  & 5 \textcolor{red}{(+1)}                 \\ \hline
                        Totale        & 3   & 3,5 & 3,5 & 5,5 & -  & 14 \textcolor{red}{(+2)} &                    \\ \hline
                        \end{tabular}
                    \caption{Consuntivo orario per ruolo nello Sprint 1}
                \end{table}

                % TODO modificare colonna "Ore Rimanenti"
                \paragraph{Prospetto Economico}\mbox{}\\
                \begin{table}[H]
                    \centering
                    \begin{tabular}{|c|c|c|c|}
                    \hline
                    \textbf{Ruolo}  & \textbf{Ore}  & \textbf{Costo} & \textbf{Ore rimanenti} \\ \hline
                    Responsabile    & 3             & €90            & 51                     \\ \hline
                    Amministratore  & 3,5           & €70            & 60,5                   \\ \hline
                    Analista        & 3,5           & €87,5          & 61,5                   \\ \hline
                    Progettista     & 5,5           & €137,5         & 99,5                   \\ \hline
                    Programmatore   & -            & -           & 184                    \\ \hline
                    Verificatore    & 14            & €210           & 179                    \\ \hline
                    \textbf{Totale} & \textbf{29,5} & \textbf{\euro595}   & \textbf{635,5}         \\ \hline
                    \end{tabular}
                    \caption{Prospetto economico e ore rimanenti}
                \end{table}

                \paragraph{Rischi effettivamente occorsi}\mbox{}\\
                RT1 - Difficoltà nell’uso di nuove tecnologie \\
                L’impatto del rischio occorso è stato minimo, visto anche la corretta attuazione della relativa misura di mitigazione preventivata, ovvero attraverso lo studio delle nuove tecnologie in modo da avere una preparazione adeguata durante lo svolgimento delle \textit{issue} previste nello sprint. \\
                \underline{Possibili migliorie:} Nessuna

                \paragraph{Retrospettiva}
                \paragraph{Risultati Ottenuti}
                    \begin{itemize}
                        \item Approvazione della candidatura per \href{https://code7crusaders.github.io/docs/PB/documentazione_interna/glossario.html#llm-large-language-model}{\textbf{LLM\textsuperscript{G}}: Assistente virtuale}.
                        \item Acquisizione di competenze sull'uso di \emph{ProjectBoard} e \emph{Roadmap} su GitHub.
                        \item Introduzione del processo di caricamento e compilazione automatizzata di file LaTeX tramite \href{https://code7crusaders.github.io/docs/PB/documentazione_interna/glossario.html#github-actions}{\emph{GitHub Actions}\textsuperscript{G}}.
                        \item Creazione della prima versione del sito statico per la documentazione.
                        \item Miglioramento dei template LaTeX per agevolare la stesura della documentazione.
                    \end{itemize}

                    \paragraph{Difficoltà Riscontrate}
                    \begin{itemize}
                        \item Pianificazione degli incontri settimanali.
                        \item Automazione del workflow per la compilazione dei file LaTeX.
                        \item Deployment automatizzato dei file sul sito web.
                    \end{itemize}
            %%%%%%%%%%%%%%%%%%%%%%%%%%%%%%%%%%%%%%%%%%%%%%%%%%%%%%%%%%%%%%
            \subsubsection{Sprint 2}
                \paragraph{Pianificazione}
                \begin{itemize}
                    \item inizio: 11 Novembre 2024
                    \item fine: 17 Novembre 2024
                \end{itemize}
                \paragraph{Sprint Planning}
                    \begin{itemize}
                        \item Obiettivo Sprint: Aggiornare il sito del gruppo, sviluppare il glossario e le norme di progetto, e condividere risorse utili tra i membri.
                        \item Attività Pianificate:
                            \begin{itemize}
                                \item Aggiornamento del sito con documenti fruibili.
                                \item Sviluppo del glossario interattivo e delle norme di progetto.
                                \item Approfondimento delle tecnologie da utilizzare per il progetto.
                            \end{itemize}
                        \item Collaborazioni:
                            \begin{itemize}
                                \item Incontro Zoom con Ergon per definire il tech stack, i requisiti, e il target del progetto.
                            \end{itemize}
                    \end{itemize}

                    \paragraph{Rischi attesi}\mbox{}\\
                    I componenti del Team ritengono che per questo Sprint siano possibili i seguenti rischi:
                    \begin{itemize}
                        \item RT1 - Difficoltà nell’uso di nuove tecnologie
                        \item RI1 - Scarsa collaborazione da parte di uno o più membri del team
                        \item RI2 - Impegni personali e universitari
                    \end{itemize}

                % TODO sposta i ruoli in base allo sprint
                \paragraph{Preventivo}\mbox{}\\
                \begin{table}[H]
                    \centering
                    \begin{tabular}{|c|c|c|c|c|c|c|c|}
                    \hline
                                  & Re  & Am  & An  & Pj  & Pg  & Ve  & Totale per persona \\ \hline
                    Cotti Cottini & -   & -   & -   & -   & -   & 4   & 4                  \\ \hline
                    Di Pietro     & 3   & -   & -   & -   & -   & -   & 3                  \\ \hline
                    Diviesti      & -   & -   & 3,5 & -   & -   & -   & 3,5                \\ \hline
                    Lapenna       & -   & -   & -   & 5,5 & -   & -   & 5,5                \\ \hline
                    Pan           & -   & -   & -   & -   & -   & 4   & 4                  \\ \hline
                    Pinarello     & -   & -   & -   & -   & -  & 4   & 4                 \\ \hline
                    Rizzolo       & -   & 3,5 & -   & -   & -   & -   & 3,5                \\ \hline
                    Totale        & 3   & 3,5 & 3,5 & 5,5 & -  & 12  &                    \\ \hline
                    \end{tabular}
                    \caption{Preventivo orario per ruolo nello Sprint 2}
                \end{table}


                \begin{figure}[H]
                    \centering
                    \begin{tikzpicture}
                        \pie[polar]{11/Responsabile, 13/Amministratore, 13/Analista, 20/Progettista, 43/Verificatore}
                    \end{tikzpicture}
                    \caption{Ripartizione in percentuale dei ruoli nello Sprint 2}
                \end{figure}


                % TODO copiala dal preventivo (per ora li facciamo uguali)
                \paragraph{Consuntivo}\mbox{}\\
                \begin{table}[H]
                    \centering
                    \begin{tabular}{|c|c|c|c|c|c|c|c|}
                        \hline
                                      & Re  & Am  & An  & Pj  & Pg  & Ve  & Totale per persona \\ \hline
                        Cotti Cottini & -   & -   & -   & -   & -   & 5 \textcolor{red}{(+1)}   & 5 \textcolor{red}{(+1)}                  \\ \hline
                        Di Pietro     & 3   & -   & -   & -   & -   & -   & 3                  \\ \hline
                        Diviesti      & -   & -   & 4,5 \textcolor{red}{(+1)} & -   & -   & -   & 4,5 \textcolor{red}{(+1)}               \\ \hline
                        Lapenna       & -   & -   & -   & 5,5 & -   & -   & 5,5                \\ \hline
                        Pan           & -   & -   & -   & -   & -   & 4   & 4                  \\ \hline
                        Pinarello     & -   & -   & -   & -   & -  & 4   & 4                 \\ \hline
                        Rizzolo       & -   & 3,5 & -   & -   & -   & -   & 3,5                \\ \hline
                        Totale        & 3   & 3,5 & 4,5 \textcolor{red}{(+1)} & 5,5 & -  & 13 \textcolor{red}{(+1)}  &                    \\ \hline
                        \end{tabular}
                    \caption{Consuntivo orario per ruolo nello Sprint 2}
                \end{table}

                % TODO modificare colonna "Ore Rimanenti"
                \paragraph{Prospetto Economico}\mbox{}\\
                \begin{table}[H]
                    \centering
                    \begin{tabular}{|c|c|c|c|}
                    \hline
                    \textbf{Ruolo}  & \textbf{Ore}  & \textbf{Costo} & \textbf{Ore rimanenti} \\ \hline
                    Responsabile    & 3             & €90            & 48                     \\ \hline
                    Amministratore  & 3,5           & €70            & 57                   \\ \hline
                    Analista        & 4,5           & €87,5          & 58                   \\ \hline
                    Progettista     & 5,5           & €137,5         & 94                   \\ \hline
                    Programmatore   & -            & -           & 184                    \\ \hline
                    Verificatore    & 13            & €150           & 166                    \\ \hline
                    \textbf{Totale} & \textbf{29,5} & \textbf{\euro535}   & \textbf{607}         \\ \hline
                    \end{tabular}
                    \caption{Prospetto economico e ore rimanenti}
                \end{table}


                \paragraph{Rischi effettivamente occorsi}\mbox{}\\
                Nessuno

                \paragraph{Retrospettiva}
                \begin{itemize}
                    \item Risultati Ottenuti:
                        \begin{itemize}
                            \item Sito del gruppo aggiornato per una migliore fruibilità dei documenti.
                            \item \href{https://code7crusaders.github.io/docs/PB/documentazione_interna/glossario.html#glossario}{Glossario\textsuperscript{G}} e \href{https://code7crusaders.github.io/docs/PB/documentazione_interna/glossario.html#norme-di-progetto}{Norme di Progetto\textsuperscript{G}} avviati.
                            \item Maggiore comprensione del progetto grazie alla condivisione delle risorse.
                        \end{itemize}
                    \item Feedback:
                        \begin{itemize}
                            \item Necessità di migliorare il passaggio da LaTeX a Markdown.
                            \item Richiesta di chiarimenti dall'azienda su alcuni use case.
                        \end{itemize}
                    \item Questioni Aperte:
                        \begin{itemize}
                            \item Come integrare automazioni e script per migliorare la valutazione del progetto.
                            \item Verifica del livello di accessibilità richiesto per il glossario.
                        \end{itemize}
                \end{itemize}
            %%%%%%%%%%%%%%%%%%%%%%%%%%%%%%%%%%%%%%%%%%%%%%%%%%%%%%%%%%%%%%
            \subsubsection{Sprint 3}
                \paragraph{Pianificazione}
                \begin{itemize}
                    \item inizio: 18 Novembre 2024
                    \item fine: 24 Novembre 2024
                \end{itemize}
                \begin{itemize}
                    \item \textbf{Obiettivi dello sprint}:
                        \begin{itemize}
                            \item Migliorare la visualizzazione dei documenti sul sito del gruppo.
                            \item Completare la stesura delle norme di progetto e dell'analisi dei requisiti.
                            \item Definire una turnazione stabile dei ruoli.
                        \end{itemize}
                    \item \textbf{Attività pianificate}:
                        \begin{itemize}
                            \item Proseguire nello sviluppo dei documenti.
                            \item Continuare l'apprendimento delle tecnologie necessarie per il progetto.
                            \item Contattare l'azienda per le specifiche hardware.
                        \end{itemize}
                    \item \textbf{Rischi e criticità}:
                        \begin{itemize}
                            \item Identificare i casi d'uso del progetto.
                            \item Valutare l'efficienza di una ricerca manuale o automatica per il glossario.
                        \end{itemize}
                \end{itemize}

                \paragraph{Rischi attesi}\mbox{}\\
                I componenti del Team ritengono che per questo Sprint siano possibili i seguenti rischi:
                \begin{itemize}
                    \item RC3 - Mancato confronto con l’azienda Proponente
                    \item RO1 - Incomprensioni nella pianificazione delle attività
                \end{itemize}

                % TODO sposta i ruoli in base allo sprint
                \paragraph{Preventivo}\mbox{}\\
                \begin{table}[H]
                    \centering
                    \begin{tabular}{|c|c|c|c|c|c|c|c|}
                    \hline
                                    & Re  & Am  & An  & Pj  & Pg  & Ve  & Totale per persona \\ \hline
                    Cotti Cottini & -   & -   & 3,5 & -   & -   & -   & 3,5                \\ \hline
                    Di Pietro     & -   & -   & -   & 5,5 & -   & -   & 5,5                \\ \hline
                    Diviesti      & -   & -   & -   & -   & -   & 5   & 5                  \\ \hline
                    Lapenna       & -   & -   & -   & -   & -   & 5   & 5                  \\ \hline
                    Pan           & -   & -   & -   & -   & 7  & -   & 7                 \\ \hline
                    Pinarello     & 3   & -   & -   & -   & -   & -   & 3                  \\ \hline
                    Rizzolo       & -   & 3,5 & -   & -   & -   & -   & 3,5                \\ \hline
                    Totale        & 3   & 3,5 & 3,5 & 5,5 & 7  & 10  &                    \\ \hline
                    \end{tabular}
                    \caption{Preventivo orario per ruolo nello Sprint 3}
                \end{table}


                \begin{figure}[H]
                    \centering
                    \begin{tikzpicture}
                        \pie[polar]{9/Responsabile, 10/Amministratore, 10/Analista, 16/Progettista, 21/Programmatore, 30/Verificatore}
                    \end{tikzpicture}
                    \caption{Ripartizione in percentuale dei ruoli nello Sprint 3}
                \end{figure}


                % TODO copiala dal preventivo (per ora li facciamo uguali)
                \paragraph{Consuntivo}\mbox{}\\
                \begin{table}[H]
                    \centering
                    \begin{tabular}{|c|c|c|c|c|c|c|c|}
                    \hline
                                  & Re  & Am  & An  & Pj  & Pg  & Ve  & Totale per persona \\ \hline
                    Cotti Cottini & -   & -   & 3,5 & -   & -   & -   & 3,5                \\ \hline
                    Di Pietro     & -   & -   & -   & 5,5 & -   & -   & 5,5                \\ \hline
                    Diviesti      & -   & -   & -   & -   & -   & 5   & 5                  \\ \hline
                    Lapenna       & -   & -   & -   & -   & -   & 6 \textcolor{red}{(+1)}   & 6 \textcolor{red}{(+1)}                  \\ \hline
                    Pan           & -   & -   & -   & -   & 8 \textcolor{red}{(+1)}  & -   & 8 \textcolor{red}{(+1)}                 \\ \hline
                    Pinarello     & 3   & -   & -   & -   & -   & -   & 3                  \\ \hline
                    Rizzolo       & -   & 3,5 & -   & -   & -   & -   & 3,5                \\ \hline
                    Totale        & 3   & 3,5 & 3,5 & 5,5 & 8 \textcolor{red}{(+1)}  & 11 \textcolor{red}{(+1)}  &                    \\ \hline
                    \end{tabular}
                    \caption{Consuntivo orario per ruolo nello Sprint 3}
                \end{table}

                % TODO modificare colonna "Ore Rimanenti"
                \paragraph{Prospetto Economico}\mbox{}\\
                \begin{table}[H]
                    \centering
                    \begin{tabular}{|c|c|c|c|}
                    \hline
                    \textbf{Ruolo}  & \textbf{Ore}  & \textbf{Costo} & \textbf{Ore rimanenti} \\ \hline
                    Responsabile    & 3             & €90            & 45                     \\ \hline
                    Amministratore  & 3,5           & €70            & 53,5                   \\ \hline
                    Analista        & 3,5           & €87,5          & 54,5                   \\ \hline
                    Progettista     & 5,5           & €137,5         & 88,5                   \\ \hline
                    Programmatore   & 8            & €120           & 176                    \\ \hline
                    Verificatore    & 11            & €165           & 155                    \\ \hline
                    \textbf{Totale} & \textbf{34,5} & \textbf{\euro670}   & \textbf{572,5}         \\ \hline
                    \end{tabular}
                    \caption{Prospetto economico e ore rimanenti}
                \end{table}


                \paragraph{Rischi effettivamente occorsi}\mbox{}\\
                RO1 - Incomprensioni nella pianificazione delle attività \\
                Vista la quantità di documentazione da aggiornare costantemente e lo svolgimento in parallelo di altre attività volte a portare avanti il progetto, il Team ha bisogno di un metodo definitivo per organizzare i lavori all’interno degli sprint. \\
                \underline{Possibili migliorie:} uno degli obbiettivi stabiliti in questo sprint, ovvero la creazione di una tabella che definisca la turnazione dei ruoli all’interno del Team

                \paragraph{Retrospettiva}
                \begin{itemize}
                    \item \textbf{Risultati raggiunti}:
                        \begin{itemize}
                            \item Aggiornamento del sito per una migliore visualizzazione dei documenti.
                            \item Condivisione di risorse utili tra i membri del team.
                            \item Buon avanzamento nella stesura delle norme di progetto e dell'analisi dei requisiti.
                            \item Turnazione dei ruoli definita in modo definitivo.
                        \end{itemize}
                    \item \textbf{Criticità risolte}:
                        \begin{itemize}
                            \item Discussione preliminare sui casi d'uso del progetto.
                        \end{itemize}
                    \item \textbf{Punti ancora da affrontare}:
                        \begin{itemize}
                            \item Creazione di una bozza di progetto per aiutare nella definizione dei requisiti hardware e software.
                        \end{itemize}
                \end{itemize}
            %%%%%%%%%%%%%%%%%%%%%%%%%%%%%%%%%%%%%%%%%%%%%%%%%%%%%%%%%%%%%%
            \subsubsection{Sprint 4}
                \paragraph{Pianificazione}
                \begin{itemize}
                    \item inizio: 25 Novembre 2024
                    \item fine: 1 Dicembre 2024
                \end{itemize}
                \begin{itemize}
                    \item Obiettivi principali:
                    \begin{itemize}
                        \item Completare la stesura del \href{https://code7crusaders.github.io/docs/PB/documentazione_interna/glossario.html#piano-di-progetto}{Piano di Progetto\textsuperscript{G}}.
                        \item Rifinire l'\href{https://code7crusaders.github.io/docs/PB/documentazione_interna/glossario.html#analisi-dei-requisiti}{Analisi dei Requisiti\textsuperscript{G}} e contattare l'azienda per feedback sui Casi d'uso.
                        \item Avanzare nello studio del framework Bloom e testare le API di ChatGPT.
                        \end{itemize}
                    \item Task assegnati:
                    \begin{itemize}
                        \item Continuare la stesura del Piano di Progetto e l'\href{https://code7crusaders.github.io/docs/PB/documentazione_interna/glossario.html#analisi-dei-requisiti}{Analisi dei Requisiti\textsuperscript{G}} \href{https://code7crusaders.github.io/docs/PB/documentazione_interna/glossario.html#piano-di-progetto}{Piano di Progetto\textsuperscript{G}} e l'Analisi dei Requisiti.
                        \item Script Python per il \href{https://code7crusaders.github.io/docs/PB/documentazione_interna/glossario.html#glossario}{Glossario\textsuperscript{G}}: debugging e miglioramenti.
                        \item Pianificare e fissare un incontro con il prof. Cardin.
                        \item Aggiornare il sito GitHub con le ultime informazioni.
                        \end{itemize}
                    \item Punti da monitorare:
                    \begin{itemize}
                        \item Progressi tecnici dello script Python per il \href{https://code7crusaders.github.io/docs/PB/documentazione_interna/glossario.html#glossario}{Glossario\textsuperscript{G}}.
                        \item Risposte dell'azienda su dataset e casi d'uso.
                        \end{itemize}
                \end{itemize}


                \paragraph{Rischi attesi}\mbox{}\\
                I componenti del Team ritengono che per questo Sprint siano possibili i seguenti rischi:
                \begin{itemize}
                    \item RT1 - Difficoltà nell’uso di nuove tecnologie
                    \item RT2 - Sviluppo di codice poco comprensibile
                    \item RO2 - Deviazione dai tempi e costi previsti
                \end{itemize}

                % TODO sposta i ruoli in base allo sprint
                \paragraph{Preventivo}\mbox{}\\
                \begin{table}[H]
                    \centering
                    \begin{tabular}{|c|c|c|c|c|c|c|c|}
                    \hline
                                  & Re  & Am  & An  & Pj  & Pg  & Ve  & Totale per persona \\ \hline
                    Cotti Cottini & -   & -   & -   & -   & -   & 5   & 5                  \\ \hline
                    Di Pietro     & -   & -   & 3,5 & -   & -   & -   & 3,5                \\ \hline
                    Diviesti      & -   & 3,5 & -   & -   & -   & -   & 3,5                \\ \hline
                    Lapenna       & -   & -   & -   & -   & 8  & -   & 8                 \\ \hline
                    Pan           & -   & -   & -   & -   & -   & 5   & 5                  \\ \hline
                    Pinarello     & -   & -   & -   & 5,5 & -   & -   & 5,5                \\ \hline
                    Rizzolo       & 3   & -   & -   & -   & -   & -   & 3                  \\ \hline
                    Totale        & 3   & 3,5 & 3,5 & 5,5 & 8  & 10  &                    \\ \hline
                    \end{tabular}
                    \caption{Preventivo orario per ruolo nello Sprint 4}
                \end{table}


                \begin{figure}[H]
                    \centering
                    \begin{tikzpicture}
                        \pie[polar]{8/Responsabile, 10/Amministratore, 10/Analista, 15/Progettista, 23/Programmatore, 28/Verificatore}
                    \end{tikzpicture}
                    \caption{Ripartizione in percentuale dei ruoli nello Sprint 4}
                \end{figure}


                % TODO copiala dal preventivo (per ora li facciamo uguali)
                \paragraph{Consuntivo}\mbox{}\\
                \begin{table}[H]
                    \centering
                    \begin{tabular}{|c|c|c|c|c|c|c|c|}
                    \hline
                                & Re  & Am  & An  & Pj  & Pg  & Ve  & Totale per persona \\ \hline
                    Cotti Cottini & -   & -   & -   & -   & -   & 6 \textcolor{red}{(+1)}  & 6 \textcolor{red}{(+1)}                 \\ \hline
                    Di Pietro     & -   & -   & 4,5 \textcolor{red}{(+1)} & -   & -   & -   & 4,5 \textcolor{red}{(+1)}               \\ \hline
                    Diviesti      & -   & 3,5 & -   & -   & -   & -   & 3,5                \\ \hline
                    Lapenna       & -   & -   & -   & -   & 8  & -   & 8                 \\ \hline
                    Pan           & -   & -   & -   & -   & -   & 5   & 5                  \\ \hline
                    Pinarello     & -   & -   & -   & 5,5 & -   & -   & 5,5                \\ \hline
                    Rizzolo       & 3   & -   & -   & -   & -   & -   & 3                  \\ \hline
                    Totale        & 3   & 3,5 & 4,5 \textcolor{red}{(+1)} & 5,5 & 8  & 11 \textcolor{red}{(+1)}  &                    \\ \hline
                    \end{tabular}
                    \caption{Consuntivo orario per ruolo nello Sprint 4}
                \end{table}

                % TODO modificare colonna "Ore Rimanenti"
                \paragraph{Prospetto Economico}\mbox{}\\
                \begin{table}[H]
                    \centering
                    \begin{tabular}{|c|c|c|c|}
                    \hline
                    \textbf{Ruolo}  & \textbf{Ore}  & \textbf{Costo} & \textbf{Ore rimanenti} \\ \hline
                    Responsabile    & 3             & €90            & 42                     \\ \hline
                    Amministratore  & 3,5           & €70            & 50                   \\ \hline
                    Analista        & 4,5           & €112,5          & 50                   \\ \hline
                    Progettista     & 5,5           & €137,5         & 83                   \\ \hline
                    Programmatore   & 8            & €120           & 168                    \\ \hline
                    Verificatore    & 11            & €165           & 144                    \\ \hline
                    \textbf{Totale} & \textbf{35,5} & \textbf{\euro695}   & \textbf{537}         \\ \hline
                    \end{tabular}
                    \caption{Prospetto economico e ore rimanenti}
                \end{table}


                \paragraph{Rischi effettivamente occorsi}\mbox{}\\
                RT1 - Difficoltà nell’uso di nuove tecnologie \\
                Il principale rischio occorso in questo sprint riguarda le tecnologie, in particolare uno script Python necessario per il Glossario. Nonostante ciò, grazie ad uno studio specifico nell’ambito del problema rilevato, anche questo rischio è stato superato senza ripercussioni gravi nelle risorse e tempistiche previste. \\
                \underline{Possibili migliorie:} Nessuna

                \paragraph{Retrospettiva}
                    \begin{itemize}
                        \item Attività completate:
                        \begin{itemize}
                            \item Documento \href{https://code7crusaders.github.io/docs/PB/documentazione_interna/glossario.html#norme-di-progetto}{Norme di Progetto\textsuperscript{G}}, inclusi processi di supporto e organizzativi.
                            \item Documento \href{https://code7crusaders.github.io/docs/PB/documentazione_interna/glossario.html#analisi-dei-requisiti}{Analisi dei Requisiti\textsuperscript{G}}, con descrizione del prodotto e casi d'uso.
                            \item Bozza dei diagrammi dei casi d'uso con relativa descrizione.
                            \item Definizione delle \textit{user-story} per i casi d'uso individuati.
                            \item Prima stesura del \href{https://code7crusaders.github.io/docs/PB/documentazione_interna/glossario.html#piano-di-progetto}{Piano di Progetto\textsuperscript{G}}.
                            \item Script Python per automatizzare il \href{https://code7crusaders.github.io/docs/PB/documentazione_interna/glossario.html#glossario}{Glossario\textsuperscript{G}}.
                        \end{itemize}
                        \item Sfide affrontate:
                        \begin{itemize}
                            \item Problemi tecnici nello script Python per il \href{https://code7crusaders.github.io/docs/PB/documentazione_interna/glossario.html#glossario}{Glossario\textsuperscript{G}}.
                            \item Individuazione e validazione dei Casi d'uso senza feedback dall'azienda.
                        \end{itemize}
                        \item Questioni aperte:
                        \begin{itemize}
                            \item Verifica della correttezza dei casi d'uso.
                            \item Dettagli sulla consegna della \href{https://code7crusaders.github.io/docs/PB/documentazione_interna/glossario.html#poc-proof-of-concept}{Proof of Concept\textsuperscript{G}} (necessità di un eseguibile?).
                        \end{itemize}
                    \end{itemize}
            %%%%%%%%%%%%%%%%%%%%%%%%%%%%%%%%%%%%%%%%%%%%%%%%%%%%%%%%%%%%%%
            \subsubsection{Sprint 5}
                \paragraph{Pianificazione}
                \begin{itemize}
                    \item inizio: 2 Dicembre 2024
                    \item fine: 8 Dicembre 2024
                \end{itemize}
                \begin{itemize}
                    \item \textbf{Obiettivi principali:}
                    \begin{itemize}
                        \item Contattare l'azienda per ottenere feedback sui casi d'uso.
                        \item Migliorare il piano di progetto per garantire allineamento con i requisiti aziendali.
                        \item Testare diversi modelli di linguaggio (\href{https://code7crusaders.github.io/docs/PB/documentazione_interna/glossario.html#llm-large-language-model}{LLM\textsuperscript{G}}) per valutarne l'idoneità.
                        \item Avviare la stesura del piano di qualifica per monitorare la qualità del progetto.
                    \end{itemize}
                    \item \textbf{Risorse necessarie:}
                    \begin{itemize}
                        \item Accesso ai dati e ai requisiti forniti dall'azienda.
                        \item Strumenti per testare e valutare i \href{https://code7crusaders.github.io/docs/PB/documentazione_interna/glossario.html#llm-large-language-model}{LLM\textsuperscript{G}}.
                        \item Linee guida e template per il piano di qualifica.
                    \end{itemize}
                    \item \textbf{Distribuzione dei compiti:}
                    \begin{itemize}
                        \item Comunicazione con l'azienda assegnata a Filippo e Francesco.
                        \item Test dei \href{https://code7crusaders.github.io/docs/PB/documentazione_interna/glossario.html#llm-large-language-model}{LLM\textsuperscript{G}} affidato a Gabriele e Eddy.
                        \item Revisione del piano di progetto gestita da Enrico e Tommaso.
                        \item Inizio del piano di qualifica a cura di Matthew.
                    \end{itemize}
                \end{itemize}


                \paragraph{Rischi attesi}\mbox{}\\
                I componenti del Team ritengono che per questo Sprint siano possibili i seguenti rischi:
                \begin{itemize}
                    \item RT1 - Difficoltà nell’uso di nuove tecnologie
                    \item RI2 - Impegni personali e universitari
                \end{itemize}

                % TODO sposta i ruoli in base allo sprint
                \paragraph{Preventivo}\mbox{}\\
                \begin{table}[H]
                    \centering
                    \begin{tabular}{|c|c|c|c|c|c|c|c|}
                    \hline
                                  & Re  & Am  & An  & Pj  & Pg  & Ve  & Totale per persona \\ \hline
                    Cotti Cottini & -   & 3,5 & -   & -   & -   & -   & 3,5                \\ \hline
                    Di Pietro     & -   & -   & -   & -   & -   & 5   & 5                  \\ \hline
                    Diviesti      & -   & -   & -   & -   & 8  & -   & 8                 \\ \hline
                    Lapenna       & -   & -   & -   & -   & -   & 5   & 5                  \\ \hline
                    Pan           & 3   & -   & -   & -   & -   & -   & 3                  \\ \hline
                    Pinarello     & -   & -   & 2,5 & -   & -   & -   & 2,5                \\ \hline
                    Rizzolo       & -   & -   & -   & 5,5 & -   & -   & 5,5                \\ \hline
                    Totale        & 3   & 3,5 & 2,5 & 5,5 & 8  & 10  &                    \\ \hline
                    \end{tabular}
                    \caption{Preventivo orario per ruolo nello Sprint 5}
                \end{table}


                \begin{figure}[H]
                    \centering
                    \begin{tikzpicture}
                        \pie[polar]{9/Responsabile, 10/Amministratore, 7/Analista, 16/Progettista, 23/Programmatore, 29/Verificatore}
                    \end{tikzpicture}
                    \caption{Ripartizione in percentuale dei ruoli nello Sprint 5}
                \end{figure}


                % TODO copiala dal preventivo (per ora li facciamo uguali)
                \paragraph{Consuntivo}\mbox{}\\
                \begin{table}[H]
                    \centering
                    \begin{tabular}{|c|c|c|c|c|c|c|c|}
                    \hline
                                    & Re  & Am  & An  & Pj  & Pg  & Ve  & Totale per persona \\ \hline
                    Cotti Cottini & -   & 3,5 & -   & -   & -   & -   & 3,5                \\ \hline
                    Di Pietro     & -   & -   & -   & -   & -   & 5   & 5                  \\ \hline
                    Diviesti      & -   & -   & -   & -   & 10 \textcolor{red}{(+2)}  & -   & 10 \textcolor{red}{(+2)}                 \\ \hline
                    Lapenna       & -   & -   & -   & -   & -   & 6 \textcolor{red}{(+1)}  & 6 \textcolor{red}{(+1)}                  \\ \hline
                    Pan           & 3   & -   & -   & -   & -   & -   & 3                  \\ \hline
                    Pinarello     & -   & -   & 3,5 \textcolor{red}{(+1)} & -   & -   & -   & 3,5 \textcolor{red}{(+1)}               \\ \hline
                    Rizzolo       & -   & -   & -   & 5,5 & -   & -   & 5,5                \\ \hline
                    Totale        & 3   & 3,5 & 3,5 \textcolor{red}{(+1)} & 5,5 & 10 \textcolor{red}{(+2)}  & 11 \textcolor{red}{(+1)}  &                    \\ \hline
                    \end{tabular}
                    \caption{Consuntivo orario per ruolo nello Sprint 5}
                \end{table}

                % TODO modificare colonna "Ore Rimanenti"
                \paragraph{Prospetto Economico}\mbox{}\\
                \begin{table}[H]
                    \centering
                    \begin{tabular}{|c|c|c|c|}
                    \hline
                    \textbf{Ruolo}  & \textbf{Ore}  & \textbf{Costo} & \textbf{Ore rimanenti} \\ \hline
                    Responsabile    & 3             & €90            & 39                     \\ \hline
                    Amministratore  & 3,5           & €70            & 46,5                   \\ \hline
                    Analista        & 3,5           & €87,5          & 46,5                   \\ \hline
                    Progettista     & 5,5           & €137,5         & 77,5                   \\ \hline
                    Programmatore   & 10            & €150           & 158                    \\ \hline
                    Verificatore    & 11            & €165           & 133                    \\ \hline
                    \textbf{Totale} & \textbf{36,5} & \textbf{\euro700}   & \textbf{500,5}         \\ \hline
                    \end{tabular}
                    \caption{Prospetto economico e ore rimanenti}
                \end{table}


                \paragraph{Rischi effettivamente occorsi}\mbox{}\\
                RI2 - Impegni personali e universitari \\
                A causa di alcuni problemi di disponibilità all’interno del Team, la riunione prevista con l’azienda Proponente per la discussione di alcune tematiche interne ha subito dei rallentamenti e un breve posticipo causando un breve periodo di riduzione del lavoro. \\
                \underline{Possibili migliorie:} Se non è possibile avere una segnalazione anticipata sulla disponibilità dei membri del Team, è buona norma recuperare nel breve termine il tempo perso che sarebbe stato da dedicare allo svolgimento dei compiti prefissati prima dell’inizio dello sprint.

                \paragraph{Retrospettiva}
                \begin{itemize}
                    \item \textbf{Risultati ottenuti:}
                    \begin{itemize}
                        \item Feedback ricevuto dall'azienda e implementato nei casi d'uso.
                        \item Piano di progetto aggiornato con nuove milestone e dettagli.
                        \item Test completati su vari \href{https://code7crusaders.github.io/docs/PB/documentazione_interna/glossario.html#llm-large-language-model}{LLM\textsuperscript{G}}, con una shortlist di modelli idonei.
                        \item Prima versione del piano di qualifica completata.
                    \end{itemize}
                    \item \textbf{Sfide incontrate:}
                    \begin{itemize}
                        \item Comunicazione iniziale con l'azienda rallentata da problemi di disponibilità.
                        \item Difficoltà nel confronto delle performance tra \href{https://code7crusaders.github.io/docs/PB/documentazione_interna/glossario.html#llm-large-language-model}{LLM\textsuperscript{G}}.
                    \end{itemize}
                    \item \textbf{Feedback ricevuto:}
                    \begin{itemize}
                        \item Apprezzamento per l'attenzione ai dettagli nel piano di progetto.
                        \item Suggerimento di includere più metriche di valutazione per i \href{https://code7crusaders.github.io/docs/PB/documentazione_interna/glossario.html#llm-large-language-model}{LLM\textsuperscript{G}}.
                    \end{itemize}
                    \item \textbf{Prossimi passi:}
                    \begin{itemize}
                        \item Continuare la documentazione.
                        \item Collaborare con l'azienda per scegliere il \href{https://code7crusaders.github.io/docs/PB/documentazione_interna/glossario.html#llm-large-language-model}{LLM\textsuperscript{G}} definitivo.
                        \item Definire l'architettura del sistema basandosi sui requisiti consolidati.
                    \end{itemize}
                \end{itemize}
            %%%%%%%%%%%%%%%%%%%%%%%%%%%%%%%%%%%%%%%%%%%%%%%%%%%%%%%%%%%%%%
            \subsubsection{Sprint 6}
                \paragraph{Pianificazione}
                \begin{itemize}
                    \item inizio: 9 Dicembre 2024
                    \item fine: 15 Dicembre 2024
                \end{itemize}
                \begin{itemize}
                    \item Aggiornare e completare la documentazione:
                    \begin{itemize}
                        \item \href{https://code7crusaders.github.io/docs/PB/documentazione_interna/glossario.html#piano-di-progetto}{Piano di Progetto\textsuperscript{G}}
                        \item \href{https://code7crusaders.github.io/docs/PB/documentazione_interna/glossario.html#analisi-dei-requisiti}{Analisi dei Requisiti\textsuperscript{G}}
                        \item \href{https://code7crusaders.github.io/docs/PB/documentazione_interna/glossario.html#glossario}{Glossario\textsuperscript{G}}
                    \end{itemize}
                    \item Aggiornare il sito del progetto.
                    \item Scegliere il framework per la creazione dell'interfaccia grafica.
                    \item Decidere il database da utilizzare.
                    \item Contattare il prof. Cardin per fissare un colloquio.
                    \item Definire i primi test statici e di unità.
                \end{itemize}


                \paragraph{Rischi attesi}\mbox{}\\
                I componenti del Team ritengono che per questo Sprint siano possibili i seguenti rischi:
                \begin{itemize}
                    \item RT1 - Difficoltà nell’uso di nuove tecnologie
                    \item RT2 - Sviluppo di codice poco comprensibile
                \end{itemize}

                % TODO sposta i ruoli in base allo sprint
                \paragraph{Preventivo}\mbox{}\\
                \begin{table}[H]
                    \centering
                    \begin{tabular}{|c|c|c|c|c|c|c|c|}
                    \hline
                                  & Re  & Am  & An  & Pj  & Pg  & Ve  & Totale per persona \\ \hline
                    Cotti Cottini & -   & -   & -   & -   & 10  & -   & 10                 \\ \hline
                    Di Pietro     & -   & 3,5 & -   & -   & -   & -   & 3,5                \\ \hline
                    Diviesti      & -   & -   & -   & -   & -   & 5   & 5                  \\ \hline
                    Lapenna       & 3   & -   & -   & -   & -   & -   & 3                  \\ \hline
                    Pan           & -   & -   & -   & 5,5 & -   & -   & 5,5                \\ \hline
                    Pinarello     & -   & -   & -   & -   & -   & 5   & 5                  \\ \hline
                    Rizzolo       & -   & -   & 3,5 & -   & -   & -   & 3,5                \\ \hline
                    Totale        & 3   & 3,5 & 3,5 & 5,5 & 10  & 10  &                    \\ \hline
                    \end{tabular}
                    \caption{Preventivo orario per ruolo nello Sprint 6}
                \end{table}


                \begin{figure}[H]
                    \centering
                    \begin{tikzpicture}
                        \pie[polar]{8/Responsabile, 10/Amministratore, 10/Analista, 15/Progettista, 27/Programmatore, 27/Verificatore}
                    \end{tikzpicture}
                    \caption{Ripartizione in percentuale dei ruoli nello Sprint 6}
                \end{figure}


                % TODO copiala dal preventivo (per ora li facciamo uguali)
                \paragraph{Consuntivo}\mbox{}\\
                \begin{table}[H]
                    \centering
                    \begin{tabular}{|c|c|c|c|c|c|c|c|}
                    \hline
                                    & Re  & Am  & An  & Pj  & Pg  & Ve  & Totale per persona \\ \hline
                    Cotti Cottini & -   & -   & -   & -   & 10  & -   & 10                 \\ \hline
                    Di Pietro     & -   & 3,5 & -   & -   & -   & -   & 3,5                \\ \hline
                    Diviesti      & -   & -   & -   & -   & -   & 5   & 5                  \\ \hline
                    Lapenna       & 3   & -   & -   & -   & -   & -   & 3                  \\ \hline
                    Pan           & -   & -   & -   & 6,5 \textcolor{red}{(+1)} & -   & -   & 6,5 \textcolor{red}{(+1)}               \\ \hline
                    Pinarello     & -   & -   & -   & -   & -   & 5   & 5                  \\ \hline
                    Rizzolo       & -   & -   & 3,5 & -   & -   & -   & 3,5                \\ \hline
                    Totale        & 3   & 3,5 & 3,5 & 6,5 \textcolor{red}{(+1)} & 10  & 10  &                    \\ \hline
                    \end{tabular}
                    \caption{Consuntivo orario per ruolo nello Sprint 6}
                \end{table}

                % TODO modificare colonna "Ore Rimanenti"
                \paragraph{Prospetto Economico}\mbox{}\\
                \begin{table}[H]
                    \centering
                    \begin{tabular}{|c|c|c|c|}
                    \hline
                    \textbf{Ruolo}  & \textbf{Ore}  & \textbf{Costo} & \textbf{Ore rimanenti} \\ \hline
                    Responsabile    & 3             & €90            & 36                     \\ \hline
                    Amministratore  & 3,5           & €70            & 43                   \\ \hline
                    Analista        & 3,5           & €87,5          & 43                   \\ \hline
                    Progettista     & 6,5           & €162,5         & 71                   \\ \hline
                    Programmatore   & 10            & €150           & 124                    \\ \hline
                    Verificatore    & 10            & €150           & 133                    \\ \hline
                    \textbf{Totale} & \textbf{36,5} & \textbf{\euro710}   & \textbf{450}         \\ \hline
                    \end{tabular}
                    \caption{Prospetto economico e ore rimanenti}
                \end{table}


                \paragraph{Rischi effettivamente occorsi}\mbox{}\\
                RT1 - Difficoltà nell’uso di nuove tecnologie \\
                Vista la necessità in questo sprint di determinare alcune scelte tecnologiche riguardanti il database da utilizzare e il framework da usare per la creazione dell’interfaccia grafica è stato molto importante lo studio di alcune possibili alternative per lo svolgimento di questi compiti selezionando ciò che abbiamo ritenuto fondamentale per le necessità richieste dal nostro prodotto. \\
                \underline{Possibili migliorie:} Nessuna

                \paragraph{Retrospettiva}
                \begin{itemize}
                    \item Attività completate:
                    \begin{itemize}
                        \item Modifiche nel documento \href{https://code7crusaders.github.io/docs/PB/documentazione_interna/glossario.html#piano-di-progetto}{Piano di Progetto\textsuperscript{G}}.
                        \item Sviluppo delle metriche di monitoraggio del progetto con relativi grafici (\href{https://code7crusaders.github.io/docs/PB/documentazione_interna/glossario.html#piano-di-qualifica}{Piano di Qualifica\textsuperscript{G}}).
                        \item Modifica e approvazione dei casi d'uso da parte dell'Azienda.
                        \item Test in locale di vari modelli \href{https://code7crusaders.github.io/docs/PB/documentazione_interna/glossario.html#llm-large-language-model}{LLM\textsuperscript{G}} per scegliere il più adeguato.
                        \item Riunione con l'Azienda e scelta definitiva del modello (gpt-o4-mini di OpenAI).
                    \end{itemize}
                    \item Sfide incontrate:
                    \begin{itemize}
                        \item Comprendere il funzionamento di \href{https://code7crusaders.github.io/docs/PB/documentazione_interna/glossario.html#langchain}{LangChain\textsuperscript{G}}.
                        \item Analisi per selezionare il miglior modello \href{https://code7crusaders.github.io/docs/PB/documentazione_interna/glossario.html#llm-large-language-model}{LLM\textsuperscript{G}} in termini di qualità-prezzo.
                    \end{itemize}
                \end{itemize}
            %%%%%%%%%%%%%%%%%%%%%%%%%%%%%%%%%%%%%%%%%%%%%%%%%%%%%%%%%%%%%%
            \subsubsection{Sprint 7}
                \paragraph{Pianificazione}
                \begin{itemize}
                    \item inizio: 16 Dicembre 2024
                    \item fine: 22 Dicembre 2024
                \end{itemize}
                Lo sprint ha avuto come obiettivi principali:
                \begin{itemize}
                    \item Stesura e revisione delle metriche di qualità;
                    \item Aggiornamento delle norme di progetto;
                    \item Definizione dei processi di base e delle metriche per i processi primari di fornitura;
                    \item Stesura e versionamento del verbale del 19/12/2024;
                    \item Integrazione di documenti firmati dall'azienda.
                \end{itemize}


                \paragraph{Rischi attesi}\mbox{}\\
                I componenti del Team ritengono che per questo Sprint siano possibili i seguenti rischi:
                \begin{itemize}
                    \item RO1 - Incomprensioni nella pianificazione delle attività
                    \item RO2 - Deviazione dai tempi e costi previsti
                \end{itemize}

                % TODO sposta i ruoli in base allo sprint
                \paragraph{Preventivo}\mbox{}\\
                \begin{table}[H]
                    \centering
                    \begin{tabular}{|c|c|c|c|c|c|c|c|}
                    \hline
                                  & Re  & Am  & An  & Pj  & Pg  & Ve  & Totale per persona \\ \hline
                    Cotti Cottini & -   & -   & -   & -   & -   & 5   & 5                  \\ \hline
                    Di Pietro     & -   & -   & -   & -   & 10  & -   & 10                 \\ \hline
                    Diviesti      & 3   & -   & -   & -   & -   & -   & 3                  \\ \hline
                    Lapenna       & -   & -   & -   & 5,5 & -   & -   & 5,5                \\ \hline
                    Pan           & -   & -   & 3,5 & -   & -   & -   & 3,5                \\ \hline
                    Pinarello     & -   & 3,5 & -   & -   & -   & -   & 3,5                \\ \hline
                    Rizzolo       & -   & -   & -   & -   & -   & 5   & 5                  \\ \hline
                    Totale        & 3   & 3,5 & 3,5 & 5,5 & 10  & 10  &                    \\ \hline
                    \end{tabular}
                    \caption{Preventivo orario per ruolo nello Sprint 7}
                \end{table}


                \begin{figure}[H]
                    \centering
                    \begin{tikzpicture}
                        \pie[polar]{8/Responsabile, 10/Amministratore, 10/Analista, 15/Progettista, 28/Programmatore, 28/Verificatore}
                    \end{tikzpicture}
                    \caption{Ripartizione in percentuale dei ruoli nello Sprint 7}
                \end{figure}


                % TODO copiala dal preventivo (per ora li facciamo uguali)
                \paragraph{Consuntivo}\mbox{}\\
                \begin{table}[H]
                    \centering
                    \begin{tabular}{|c|c|c|c|c|c|c|c|}
                    \hline
                                    & Re  & Am  & An  & Pj  & Pg  & Ve  & Totale per persona \\ \hline
                    Cotti Cottini & -   & -   & -   & -   & -   & 7 \textcolor{red}{(+2)}   & 7 \textcolor{red}{(+2)}                  \\ \hline
                    Di Pietro     & -   & -   & -   & -   & 10  & -   & 10                 \\ \hline
                    Diviesti      & 3   & -   & -   & -   & -   & -   & 3                  \\ \hline
                    Lapenna       & -   & -   & -   & 6,5 \textcolor{red}{(+1)} & -   & -   & 6,5 \textcolor{red}{(+1)}               \\ \hline
                    Pan           & -   & -   & 3,5 & -   & -   & -   & 3,5                \\ \hline
                    Pinarello     & -   & 3,5 & -   & -   & -   & -   & 3,5                \\ \hline
                    Rizzolo       & -   & -   & -   & -   & -   & 5   & 5                  \\ \hline
                    Totale        & 3   & 3,5 & 3,5 & 6,5 \textcolor{red}{(+1)} & 10  & 12 \textcolor{red}{(+2)}  &                    \\ \hline
                    \end{tabular}
                    \caption{Consuntivo orario per ruolo nello Sprint 7}
                \end{table}

                % TODO modificare colonna "Ore Rimanenti"
                \paragraph{Prospetto Economico}\mbox{}\\
                \begin{table}[H]
                    \centering
                    \begin{tabular}{|c|c|c|c|}
                    \hline
                    \textbf{Ruolo}  & \textbf{Ore}  & \textbf{Costo} & \textbf{Ore rimanenti} \\ \hline
                    Responsabile    & 3             & €90            & 33                     \\ \hline
                    Amministratore  & 3,5           & €70            & 39,5                   \\ \hline
                    Analista        & 3,5           & €87,5          & 39,5                   \\ \hline
                    Progettista     & 6,5           & €162,5         & 64,5                   \\ \hline
                    Programmatore   & 10            & €150           & 114                    \\ \hline
                    Verificatore    & 12            & €180           & 121                    \\ \hline
                    \textbf{Totale} & \textbf{38,5} & \textbf{\euro740}   & \textbf{411,5}         \\ \hline
                    \end{tabular}
                    \caption{Prospetto economico e ore rimanenti}
                \end{table}


                \paragraph{Rischi effettivamente occorsi}\mbox{}\\
                RO2 - Deviazione dai tempi e costi previsti \\
                Questo rischio si è verificato a causa di alcuni problemi e piccoli errori in alcuni documenti che trattano argomenti tra cui le metriche di qualità. Per via di ciò, il Team ha dovuto impiegare tempo extra per la correzione e la revisione di questi documenti in modo da avere una documentazione al passo con lo standard qualitativo previsto. \\
                \underline{Possibili migliorie:} Cercare di evitare errori ed effettuare una revisione nel breve termine dei documenti che hanno subito modifiche importanti

                \paragraph{Retrospettiva}
                Durante lo sprint, il team è riuscito a completare gli obiettivi prefissati:
                \begin{itemize}
                    \item Sono state corrette e aggiornate le metriche di qualità;
                    \item È stato effettuato un merge delle modifiche alle norme di progetto;
                    \item È stato definito uno standard per la qualità;
                    \item Il verbale del 19/12/2024 è stato finalizzato e versionato correttamente;
                    \item Sono stati integrati i documenti firmati dall'azienda.
                \end{itemize}
                
                Alcune difficoltà incontrate durante lo sprint:
                \begin{itemize}
                    \item Problemi con la gestione delle metriche di qualità, che hanno richiesto una correzione successiva;
                    \item Necessità di ripristinare una versione precedente dei documenti a causa di errori;
                    \item Tempo extra dedicato alla rifinitura e revisione dei documenti per garantire la qualità.
                \end{itemize}
                
                Per migliorare l'efficienza nei prossimi sprint, il team ha deciso di:
                \begin{itemize}
                    \item Pianificare con maggiore precisione le revisioni, evitando modifiche last-minute;
                    \item Adottare una gestione più strutturata delle metriche di qualità;
                    \item Ottimizzare il processo di merge per ridurre i conflitti tra le modifiche dei documenti.
                \end{itemize}
            %%%%%%%%%%%%%%%%%%%%%%%%%%%%%%%%%%%%%%%%%%%%%%%%%%%%%%%%%%%%%%
            \subsubsection{Sprint 8}
                \paragraph{Pianificazione}
                \begin{itemize}
                    \item inizio: 23 Dicembre 2024
                    \item fine: 29 Dicembre 2024
                \end{itemize}
                L'obiettivo principale di questo sprint è stato la creazione della documentazione essenziale per il progetto, inclusi i seguenti elementi:
                \begin{itemize}
                    \item Creazione del template base per la gestione dello sprint.
                    \item Redazione del piano di progetto (versione 0.3).
                    \item Analisi delle tecnologie per il frontend e scelta del framework.
                    \item Correzione e miglioramento della documentazione di analisi frontend.
                \end{itemize}
                
                Le attività previste per questo sprint includono:
                \begin{itemize}
                    \item Creazione di un documento per la scelta del framework frontend.
                    \item Correzione di errori nel documento di analisi frontend.
                    \item Creazione del template base per lo sprint.
                    \item Redazione della versione 0.3 del piano di progetto.
                \end{itemize}


                \paragraph{Rischi attesi}\mbox{}\\
                I componenti del Team ritengono che per questo Sprint siano possibili i seguenti rischi:
                \begin{itemize}
                    \item RT1 - Difficoltà nell’uso di nuove tecnologie
                    \item RO2 - Deviazione dai tempi e costi previsti
                    \item RI2 - Impegni personali e universitari
                \end{itemize}

                % TODO sposta i ruoli in base allo sprint
                \paragraph{Preventivo}\mbox{}\\
                \begin{table}[H]
                    \centering
                    \begin{tabular}{|c|c|c|c|c|c|c|c|}
                    \hline
                                  & Re  & Am  & An  & Pj  & Pg  & Ve  & Totale per persona \\ \hline
                    Cotti Cottini & 3   & -   & -   & -   & -   & -   & 3                  \\ \hline
                    Di Pietro     & -   & -   & -   & -   & -   & 5   & 5                  \\ \hline
                    Diviesti      & -   & -   & -   & 5,5 & -   & -   & 5,5                \\ \hline
                    Lapenna       & -   & -   & 3,5 & -   & -   & -   & 3,5                \\ \hline
                    Pan           & -   & -   & -   & -   & -   & 5   & 5                  \\ \hline
                    Pinarello     & -   & -   & -   & -   & 8  & -   & 8                 \\ \hline
                    Rizzolo       & -   & 3,5 & -   & -   & -   & -   & 3,5                \\ \hline
                    Totale        & 3   & 3,5 & 3,5 & 5,5 & 8  & 10  &                    \\ \hline
                    \end{tabular}
                    \caption{Preventivo orario per ruolo nello Sprint 8}
                \end{table}
                

                \begin{figure}[H]
                    \centering
                    \begin{tikzpicture}
                        \pie[polar]{9/Responsabile, 10/Amministratore, 10/Analista, 16/Progettista, 23/Programmatore, 29/Verificatore}
                    \end{tikzpicture}
                    \caption{Ripartizione in percentuale dei ruoli nello Sprint 8}
                \end{figure}

            
                % TODO copiala dal preventivo (per ora li facciamo uguali)
                \paragraph{Consuntivo}\mbox{}\\
                \begin{table}[H]
                    \centering
                    \begin{tabular}{|c|c|c|c|c|c|c|c|}
                    \hline
                                    & Re  & Am  & An  & Pj  & Pg  & Ve  & Totale per persona \\ \hline
                    Cotti Cottini & 3   & -   & -   & -   & -   & -   & 3                  \\ \hline
                    Di Pietro     & -   & -   & -   & -   & -   & 7 \textcolor{red}{(+2)}   & 7 \textcolor{red}{(+2)}                  \\ \hline
                    Diviesti      & -   & -   & -   & 6,5 \textcolor{red}{(+1)} & -   & -   & 6,5 \textcolor{red}{(+1)}                \\ \hline
                    Lapenna       & -   & -   & 3,5 & -   & -   & -   & 3,5                \\ \hline
                    Pan           & -   & -   & -   & -   & -   & 5   & 5                  \\ \hline
                    Pinarello     & -   & -   & -   & -   & 8  & -   & 8                 \\ \hline
                    Rizzolo       & -   & 3,5 & -   & -   & -   & -   & 3,5                \\ \hline
                    Totale        & 3   & 3,5 & 3,5 & 6,5 \textcolor{red}{(+1)} & 8  & 12 \textcolor{red}{(+2)}  &                    \\ \hline
                    \end{tabular}
                    \caption{Consuntivo orario per ruolo nello Sprint 8}
                \end{table}

                % TODO modificare colonna "Ore Rimanenti"
                \paragraph{Prospetto Economico}\mbox{}\\
                \begin{table}[H]
                    \centering
                    \begin{tabular}{|c|c|c|c|}
                    \hline
                    \textbf{Ruolo}  & \textbf{Ore}  & \textbf{Costo} & \textbf{Ore rimanenti} \\ \hline
                    Responsabile    & 3             & €90            & 30                     \\ \hline
                    Amministratore  & 3,5           & €70            & 36                   \\ \hline
                    Analista        & 3,5           & €87,5          & 36                   \\ \hline
                    Progettista     & 6,5           & €162,5         & 58                   \\ \hline
                    Programmatore   & 8            & €120           & 106                    \\ \hline
                    Verificatore    & 12            & €180           & 109                    \\ \hline
                    \textbf{Totale} & \textbf{36,5} & \textbf{\euro710}   & \textbf{375}         \\ \hline
                    \end{tabular}
                    \caption{Prospetto economico e ore rimanenti}
                \end{table}


                \paragraph{Rischi effettivamente occorsi}\mbox{}\\
                Nessuno

                \paragraph{Retrospettiva}
                Durante lo sprint sono stati completati i seguenti task:
                \begin{itemize}
                    \item È stato creato un template base per la gestione dello sprint.
                    \item È stata completata la versione 0.3 del piano di progetto.
                    \item È stato corretto e migliorato il documento di analisi frontend.
                    \item È stato creato un documento di valutazione per la scelta del framework frontend.
                \end{itemize}
                
                Feedback e Miglioramenti:
                \begin{itemize}
                    \item Il template base per lo sprint potrebbe essere ulteriormente ottimizzato per facilitare l'inserimento delle informazioni.
                    \item La scelta del framework frontend potrebbe richiedere ulteriori approfondimenti e prove pratiche.
                    \item La documentazione di analisi frontend è stata migliorata, ma potrebbe necessitare di ulteriori revisioni per garantire chiarezza e coerenza.
                \end{itemize}
                
                Prossimi Passi:
                \begin{itemize}
                    \item Refinement del template base per lo sprint.
                    \item Validazione finale della scelta del framework frontend.
                    \item Revisione continua della documentazione per mantenerla aggiornata e coerente con il progetto.
                \end{itemize}
            %%%%%%%%%%%%%%%%%%%%%%%%%%%%%%%%%%%%%%%%%%%%%%%%%%%%%%%%%%%%%%
            \subsubsection{Sprint 9}
                \paragraph{Pianificazione}
                \begin{itemize}
                    \item inizio: 30 Dicembre 2024
                    \item fine: 5 Gennaio 2025
                \end{itemize}
                Durante questo sprint, il team si è concentrato sulla compilazione e revisione della documentazione, con particolare attenzione al \textbf{PDP} (Project Development Plan) e alla documentazione dello sprint.

                Le attività previste per lo sprint includevano:
                \begin{itemize}
                    \item Compilazione della documentazione \textbf{PDP}
                    \item Aggiornamenti minori e revisione documenti
                    \item Compilazione della documentazione \textbf{Sprint 1}
                    \item Automazione della compilazione PDF tramite GitHub Actions
                \end{itemize}


                \paragraph{Rischi attesi}\mbox{}\\
                I componenti del Team ritengono che per questo Sprint siano possibili i seguenti rischi:
                \begin{itemize}
                    \item RT1 - Difficoltà nell’uso di nuove tecnologie
                    \item RC2 - Problemi di comunicazione all’interno del gruppo
                    \item RC3 - Mancato confronto con l’azienda Proponente
                \end{itemize}

                % TODO sposta i ruoli in base allo sprint
                \paragraph{Preventivo}\mbox{}\\
                \begin{table}[H]
                    \centering
                    \begin{tabular}{|c|c|c|c|c|c|c|c|}
                    \hline
                                  & Re  & Am  & An  & Pj  & Pg  & Ve  & Totale per persona \\ \hline
                    Cotti Cottini & -   & -   & -   & 5,5 & -   & -   & 5,5                \\ \hline
                    Di Pietro     & 3   & -   & -   & -   & -   & -   & 3                  \\ \hline
                    Diviesti      & -   & -   & 3,5 & -   & -   & -   & 3,5                \\ \hline
                    Lapenna       & -   & -   & -   & -   & -   & 5   & 5                  \\ \hline
                    Pan           & -   & 3,5 & -   & -   & -   & -   & 3,5                \\ \hline
                    Pinarello     & -   & -   & -   & -   & -   & 5   & 5                  \\ \hline
                    Rizzolo       & -   & -   & -   & -   & 10  & -   & 10                 \\ \hline
                    Totale        & 3   & 3,5 & 3,5 & 5,5 & 10  & 10  &                    \\ \hline
                    \end{tabular}
                    \caption{Preventivo orario per ruolo nello Sprint 9}
                \end{table}


                \begin{figure}[H]
                    \centering
                    \begin{tikzpicture}
                        \pie[polar]{8/Responsabile, 10/Amministratore, 10/Analista, 15/Progettista, 28/Programmatore, 28/Verificatore}
                    \end{tikzpicture}
                    \caption{Ripartizione in percentuale dei ruoli nello Sprint 9}
                \end{figure}


                % TODO copiala dal preventivo (per ora li facciamo uguali)
                \paragraph{Consuntivo}\mbox{}\\
                \begin{table}[H]
                    \centering
                    \begin{tabular}{|c|c|c|c|c|c|c|c|}
                    \hline
                                    & Re  & Am  & An  & Pj  & Pg  & Ve  & Totale per persona \\ \hline
                    Cotti Cottini & -   & -   & -   & 5,5 & -   & -   & 5,5                \\ \hline
                    Di Pietro     & 4 \textcolor{red}{(+1)}   & -   & -   & -   & -   & -   & 4 \textcolor{red}{(+1)}                  \\ \hline
                    Diviesti      & -   & -   & 3,5 & -   & -   & -   & 3,5                \\ \hline
                    Lapenna       & -   & -   & -   & -   & -   & 5   & 5                  \\ \hline
                    Pan           & -   & 3,5 & -   & -   & -   & -   & 3,5                \\ \hline
                    Pinarello     & -   & -   & -   & -   & -   & 6 \textcolor{red}{(+1)}   & 6 \textcolor{red}{(+1)}                  \\ \hline
                    Rizzolo       & -   & -   & -   & -   & 10  & -   & 10                 \\ \hline
                    Totale        & 4 \textcolor{red}{(+1)}   & 3,5 & 3,5 & 5,5 & 10  & 11 \textcolor{red}{(+1)}  &                    \\ \hline
                    \end{tabular}
                    \caption{Consuntivo orario per ruolo nello Sprint 9}
                \end{table}

                % TODO modificare colonna "Ore Rimanenti"
                \paragraph{Prospetto Economico}\mbox{}\\
                \begin{table}[H]
                    \centering
                    \begin{tabular}{|c|c|c|c|}
                    \hline
                    \textbf{Ruolo}  & \textbf{Ore}  & \textbf{Costo} & \textbf{Ore rimanenti} \\ \hline
                    Responsabile    & 4             & €120            & 26                     \\ \hline
                    Amministratore  & 3,5           & €70            & 32,5                   \\ \hline
                    Analista        & 3,5           & €87,5          & 32,5                   \\ \hline
                    Progettista     & 5,5           & €137,5         & 55,5                   \\ \hline
                    Programmatore   & 10            & €150           & 94                    \\ \hline
                    Verificatore    & 11            & €165           & 98                    \\ \hline
                    \textbf{Totale} & \textbf{37,5} & \textbf{\euro730}   & \textbf{338,5}         \\ \hline
                    \end{tabular}
                    \caption{Prospetto economico e ore rimanenti}
                \end{table}


                \paragraph{Rischi effettivamente occorsi}\mbox{}\\
                RT1 - Difficoltà nell’uso di nuove tecnologie \\
                L’impatto del rischio occorso è stato minimo, visto anche la corretta attuazione della relativa misura di mitigazione preventivata, ovvero attraverso lo studio delle nuove tecnologie in modo da avere una preparazione adeguata durante lo svolgimento delle \textit{issue} previste nello sprint. \\
                \underline{Possibili migliorie:} Nessuna

                \paragraph{Retrospettiva}
                Lo sprint ha portato alla realizzazione dei seguenti risultati:
                \begin{itemize}
                    \item La documentazione \textbf{PDP} è stata completata e compilata con successo il \textbf{2 gennaio 2025}.
                    \item La documentazione \textbf{Sprint 1} è stata finalizzata il \textbf{30 dicembre 2024}.
                    \item Sono stati applicati miglioramenti minori e correzioni al documento.
                    \item L'automazione della compilazione dei PDF ha funzionato correttamente, eseguendo la generazione automatica della documentazione.
                \end{itemize}
                
                Il processo di automazione ha funzionato bene, riducendo il carico manuale sulla generazione della documentazione. Tuttavia, per i prossimi sprint, si suggerisce:
                \begin{itemize}
                    \item Migliorare il sistema di naming delle versioni per facilitare il tracking delle modifiche.
                    \item Integrare un controllo qualità sui PDF generati prima della loro pubblicazione automatica.
                    \item Valutare l'aggiunta di una pipeline CI/CD più strutturata per gestire la documentazione in modo efficiente.
                \end{itemize}
            %%%%%%%%%%%%%%%%%%%%%%%%%%%%%%%%%%%%%%%%%%%%%%%%%%%%%%%%%%%%%%
            \subsubsection{Sprint 10}
                \paragraph{Pianificazione}
                \begin{itemize}
                    \item inizio: 6 Gennaio 2025
                    \item fine: 12 Gennaio 2025
                \end{itemize}
                Durante questo sprint, il team si è concentrato su diversi aspetti della documentazione di progetto e su alcune revisioni e aggiornamenti di qualità. Gli obiettivi principali erano:
                \begin{itemize}
                    \item Continuazione e aggiornamento del Piano di Progetto.
                    \item Revisione e verifica della documentazione di Norme di Progetto e Piano di Qualifica.
                    \item Correzione di errori grammaticali e miglioramenti nella qualità della documentazione.
                    \item Automazione della compilazione PDF e gestione script di compilazione.
                    \item Aggiornamento dei diari di lavoro.
                \end{itemize}


                \paragraph{Rischi attesi}\mbox{}\\
                I componenti del Team ritengono che per questo Sprint siano possibili i seguenti rischi:
                \begin{itemize}
                    \item RC2 - Problemi di comunicazione all’interno del gruppo
                    \item RC3 - Mancato confronto con l’azienda Proponente
                    \item RI1 - Scarsa collaborazione da parte di uno o più membri del team
                \end{itemize}

                \paragraph{Preventivo}\mbox{}\\
                \begin{table}[H]
                    \centering
                    \begin{tabular}{|c|c|c|c|c|c|c|c|}
                    \hline
                                  & Re  & Am  & An  & Pj  & Pg  & Ve  & Totale per persona \\ \hline
                    Cotti Cottini & -   & -   & 3,5 & -   & -   & -   & 3,5                \\ \hline
                    Di Pietro     & -   & -   & -   & 5,5 & -   & -   & 5,5                \\ \hline
                    Diviesti      & -   & -   & -   & -   & -   & 5   & 5                  \\ \hline
                    Lapenna       & -   & 3,5 & -   & -   & -   & -   & 3,5                \\ \hline
                    Pan           & -   & -   & -   & -   & 10  & -   & 10                 \\ \hline
                    Pinarello     & 3   & -   & -   & -   & -   & -   & 3                  \\ \hline
                    Rizzolo       & -   & -   & -   & -   & -   & 5   & 5                  \\ \hline
                    Totale        & 3   & 3,5 & 3,5 & 5,5 & 10  & 10  &                    \\ \hline
                    \end{tabular}
                    \caption{Preventivo orario per ruolo nello Sprint 10}
                \end{table}


                \begin{figure}[H]
                    \centering
                    \begin{tikzpicture}
                        \pie[polar]{8/Responsabile, 10/Amministratore, 10/Analista, 15/Progettista, 28/Programmatore, 28/Verificatore}
                    \end{tikzpicture}
                    \caption{Ripartizione in percentuale dei ruoli nello Sprint 10}
                \end{figure}


                % TODO copiala dal preventivo (per ora li facciamo uguali)
                \paragraph{Consuntivo}\mbox{}\\
                \begin{table}[H]
                    \centering
                    \begin{tabular}{|c|c|c|c|c|c|c|c|}
                    \hline
                                    & Re  & Am  & An  & Pj  & Pg  & Ve  & Totale per persona \\ \hline
                    Cotti Cottini & -   & -   & 3,5 & -   & -   & -   & 3,5               \\ \hline
                    Di Pietro     & -   & -   & -   & 5,5 & -   & -   & 5,5                 \\ \hline
                    Diviesti      & -   & -   & -   & -   & -   & 5   & 5               \\ \hline
                    Lapenna       & -   & 3,5 & -   & -   & -   & -   & 3,5                 \\ \hline
                    Pan           & -   & -   & -   & -   & 10  & -   & 10                \\ \hline
                    Pinarello     & 3   & -   & -   & -   & -   & -   & 3                 \\ \hline
                    Rizzolo       & -   & -   & -   & -   & -   & 5   & 5                 \\ \hline
                    Totale        & 3   & 3,5 & 3,5 & 5,5 & 10  & 10  &                     \\ \hline
                    \end{tabular}
                    \caption{Consuntivo orario per ruolo nello Sprint 10}
                \end{table}

                % TODO modificare colonna "Ore Rimanenti"
                \paragraph{Prospetto Economico}\mbox{}\\
                \begin{table}[H]
                    \centering
                    \begin{tabular}{|c|c|c|c|}
                    \hline
                    \textbf{Ruolo}  & \textbf{Ore}  & \textbf{Costo} & \textbf{Ore rimanenti} \\ \hline
                    Responsabile    & 3             & €90            & 23                     \\ \hline
                    Amministratore  & 3,5           & €70            & 29                   \\ \hline
                    Analista        & 3,5           & €87,5          & 29                   \\ \hline
                    Progettista     & 5,5           & €137,5         & 50                   \\ \hline
                    Programmatore   & 10            & €150           & 84                    \\ \hline
                    Verificatore    & 10            & €150           & 88                    \\ \hline
                    \textbf{Totale} & \textbf{35,5} & \textbf{\euro685}   & \textbf{303}         \\ \hline
                    \end{tabular}
                    \caption{Prospetto economico e ore rimanenti}
                \end{table}


                \paragraph{Rischi effettivamente occorsi}\mbox{}\\
                Nessuno

                \paragraph{Retrospettiva}
                Durante lo sprint, il team ha completato con successo le seguenti attività:
                \begin{itemize}
                    \item Il Piano di Progetto è stato aggiornato e migliorato.
                    \item Sono state effettuate revisioni su Norme di Progetto e Piano di Qualifica.
                    \item È stata implementata la correzione di errori grammaticali e migliorata la qualità della documentazione.
                    \item È stato aggiornato e verificato lo script di compilazione per la generazione automatica della documentazione in PDF.
                    \item È stato aggiunto un nuovo diario di lavoro (08/01/2025).
                    \item La documentazione di Sprint Planning e Review è stata aggiornata fino allo Sprint 7.
                \end{itemize}
                
                \textbf{Cosa ha funzionato bene:}
                \begin{itemize}
                    \item Collaborazione efficace tra i membri del team.
                    \item Processo di revisione e verifica ben strutturato.
                    \item Automazione della compilazione della documentazione, riducendo il rischio di errori manuali.
                \end{itemize}
                
                \textbf{Aree di miglioramento:}
                \begin{itemize}
                    \item Maggiore attenzione nella gestione delle versioni dei documenti.
                    \item Ottimizzazione della pipeline di compilazione per una maggiore efficienza.
                \end{itemize}
                
                \textbf{Azioni per il prossimo Sprint:}
                \begin{itemize}
                    \item Continuare la revisione della documentazione per migliorarne la qualità.
                    \item Rafforzare il processo di verifica dei documenti per ridurre le correzioni postume.
                    \item Testare ulteriormente lo script di compilazione per identificare eventuali miglioramenti.
                \end{itemize}
            %%%%%%%%%%%%%%%%%%%%%%%%%%%%%%%%%%%%%%%%%%%%%%%%%%%%%%%%%%%%%%
            \subsubsection{Sprint 11}
                \paragraph{Pianificazione}
                \begin{itemize}
                    \item inizio: 13 Gennaio 2025
                    \item fine: 19 Gennaio 2025
                \end{itemize}
                L'obiettivo principale di questo sprint è stato aggiornare e verificare la documentazione del progetto, in particolare:
                \begin{itemize}
                    \item Aggiornamento del Piano di Progetto;
                    \item Aggiornamento e verifica delle Norme di Progetto;
                    \item Revisione e aggiornamento della sezione dei casi d'uso;
                    \item Verifica e aggiornamento dell'Analisi dei Requisiti;
                    \item Implementazione della sezione di Test di Sistema v1;
                    \item Generazione automatica dei PDF aggiornati della documentazione.
                \end{itemize}


                \paragraph{Rischi attesi}\mbox{}\\
                I componenti del Team ritengono che per questo Sprint siano possibili i seguenti rischi:
                \begin{itemize}
                    \item RO1 - Incomprensioni nella pianificazione delle attività
                    \item RI2 - Impegni personali e universitari
                \end{itemize}

                \paragraph{Preventivo}\mbox{}\\
                \begin{table}[H]
                    \centering
                    \begin{tabular}{|c|c|c|c|c|c|c|c|}
                    \hline
                                  & Re  & Am  & An  & Pj  & Pg  & Ve  & Totale per persona \\ \hline
                    Cotti Cottini & -   & -   & -   & -   & -   & 5   & 5                  \\ \hline
                    Di Pietro     & -   & -   & 3,5 & -   & -   & -   & 3,5                \\ \hline
                    Diviesti      & -   & 3,5 & -   & -   & -   & -   & 3,5                \\ \hline
                    Lapenna       & -   & -   & -   & -   & 10  & -   & 10                 \\ \hline
                    Pan           & -   & -   & -   & -   & -   & 5   & 5                  \\ \hline
                    Pinarello     & -   & -   & -   & 5,5 & -   & -   & 5,5                \\ \hline
                    Rizzolo       & 3   & -   & -   & -   & -   & -   & 3                  \\ \hline
                    Totale        & 3   & 3,5 & 3,5 & 5,5 & 10  & 10  &                    \\ \hline
                    \end{tabular}
                    \caption{Preventivo orario per ruolo nello Sprint 11}
                \end{table}


                \begin{figure}[H]
                    \centering
                    \begin{tikzpicture}
                        \pie[polar]{8/Responsabile, 10/Amministratore, 10/Analista, 15/Progettista, 28/Programmatore, 28/Verificatore}
                    \end{tikzpicture}
                    \caption{Ripartizione in percentuale dei ruoli nello Sprint 11}
                \end{figure}


                % TODO copiala dal preventivo (per ora li facciamo uguali)
                \paragraph{Consuntivo}\mbox{}\\
                \begin{table}[H]
                    \centering
                    \begin{tabular}{|c|c|c|c|c|c|c|c|}
                    \hline
                                  & Re  & Am  & An  & Pj  & Pg  & Ve  & Totale per persona \\ \hline
                    Cotti Cottini & -   & -   & -   & -   & -   & 6 \textcolor{red}{(+1)}   & 6 \textcolor{red}{(+1)}                  \\ \hline
                    Di Pietro     & -   & -   & 3,5 & -   & -   & -   & 3,5                \\ \hline
                    Diviesti      & -   & 3,5 & -   & -   & -   & -   & 3,5                \\ \hline
                    Lapenna       & -   & -   & -   & -   & 9  & -   & 9                 \\ \hline
                    Pan           & -   & -   & -   & -   & -   & 6 \textcolor{red}{(+1)}   & 6 \textcolor{red}{(+1)}                  \\ \hline
                    Pinarello     & -   & -   & -   & 6,5 \textcolor{red}{(+1)} & -   & -   & 6,5 \textcolor{red}{(+1)}               \\ \hline
                    Rizzolo       & 3   & -   & -   & -   & -   & -   & 3                  \\ \hline
                    Totale        & 3   & 3,5 & 3,5 & 6,5 \textcolor{red}{(+1)} & 9  & 12 \textcolor{red}{(+1)}  &                    \\ \hline
                    \end{tabular}
                    \caption{Consuntivo orario per ruolo nello Sprint 11}
                \end{table}

                % TODO modificare colonna "Ore Rimanenti"
                \paragraph{Prospetto Economico}\mbox{}\\
                \begin{table}[H]
                    \centering
                    \begin{tabular}{|c|c|c|c|}
                    \hline
                    \textbf{Ruolo}  & \textbf{Ore}  & \textbf{Costo} & \textbf{Ore rimanenti} \\ \hline
                    Responsabile    & 3             & €90            & 20                     \\ \hline
                    Amministratore  & 3,5           & €70            & 25,5                     \\ \hline
                    Analista        & 3,5           & €87,5          & 25,5                     \\ \hline
                    Progettista     & 6,5           & €162,5         & 43,5                     \\ \hline
                    Programmatore   & 9            & €135           & 75                     \\ \hline
                    Verificatore    & 12            & €180           & 76                     \\ \hline
                    \textbf{Totale} & \textbf{37,5} & \textbf{\euro725}   & \textbf{265,5}           \\ \hline
                    \end{tabular}
                    \caption{Prospetto economico e ore rimanenti}
                \end{table}


                \paragraph{Rischi effettivamente occorsi}\mbox{}\\
                RO1 - Incomprensioni nella pianificazione delle attività \\
                Viste le tempistiche necessarie per verificare documentazione sviluppata in sprint precedenti, e visto anche lo stato di avanzamento di progetto bisognerebbe avere le idee più chiare e dei traguardi realistici e ben definiti alla fine di ogni sprint. \\
                \underline{Possibili migliorie:} all’inizio di ogni sprint vengono definiti in maniera chiara degli obbiettivi da portare a termine in modo da effettuare dei compiti superflui negli sprint successivi

                \paragraph{Retrospettiva}
                Durante lo sprint sono stati completati i seguenti task:
                \begin{itemize}
                    \item Aggiornamento del Piano di Progetto (commit: ``Update piano\_di\_progetto\_v0.5.tex'');
                    \item Aggiornamento dei rischi nel Piano di Progetto (commit: ``aggiornamento rischi'');
                    \item Verifica delle Norme di Progetto (commit: ``verificato norme di progetto'');
                    \item Revisione e aggiornamento dei casi d'uso, incluso il diagramma (commit: ``aggiornato diagramma casi d'uso'', ``Aggiornati i casi d'uso in AdR'');
                    \item Aggiunta di tabelle attuali sui requisiti (commit: ``aggiunte tabelle attuali requisiti in analisi dei requisiti'');
                    \item Implementazione della sezione Test di Sistema v1 (commit: ``PDQ: Test di Sistema v1'');
                    \item Revisione della sezione 3 delle Norme di Progetto (commit: ``Verifica sezione 3 NDP'');
                    \item Ottimizzazione della pipeline di compilazione automatica (commit: ``Update compile.sh'').
                \end{itemize}
                
                Durante la Sprint Review, abbiamo evidenziato i seguenti punti:
                \begin{itemize}
                    \item \textbf{Cosa è andato bene:} 
                    \begin{itemize}
                        \item Buona collaborazione tra i membri del team;
                        \item La documentazione è stata aggiornata con successo;
                        \item La pipeline di compilazione automatica ha funzionato correttamente.
                    \end{itemize}
                    \item \textbf{Cosa può essere migliorato:}
                    \begin{itemize}
                        \item Maggiore attenzione alla coerenza tra i documenti;
                        \item Ottimizzazione della revisione dei casi d'uso per evitare duplicazioni di lavoro;
                        \item Definizione più chiara degli obiettivi dello sprint all'inizio.
                    \end{itemize}
                \end{itemize}
            %%%%%%%%%%%%%%%%%%%%%%%%%%%%%%%%%%%%%%%%%%%%%%%%%%%%%%%%%%%%%%
            \subsubsection{Sprint 12}
                \paragraph{Pianificazione}
                \begin{itemize}
                    \item inizio: 20 Gennaio 2025
                    \item fine: 26 Gennaio 2025
                \end{itemize}
                Durante lo sprint, il team si è focalizzato sui seguenti obiettivi:
                \begin{itemize}
                    \item Aggiornamento e revisione del \textit{Piano di Qualifica} (piano\_di\_qualifica\_v0.3.tex).
                    \item Aggiornamento e revisione del \textit{Piano di Progetto} (piano\_di\_progetto\_v0.6.tex).
                    \item Aggiornamento delle \textit{Norme di Progetto} (norme\_di\_progetto.tex).
                    \item Creazione e aggiornamento della documentazione relativa ai diagrammi UML:
                    \begin{itemize}
                        \item Database (DataBase.mdj).
                        \item Diagramma dei Casi d'Uso (DiagrammaCasiUso\_v1.2.drawio).
                        \item Correzione delle condizioni extends.
                    \end{itemize}
                    \item Aggiornamento e verifica dei verbali interni e con il cliente:
                    \begin{itemize}
                        \item Verbale interno del 19/12/2024.
                        \item Verbale con il cliente del 24/12/2024.
                        \item Verbale del 29/11/2024.
                    \end{itemize}
                    \item Creazione della lettera di presentazione (lettera\_di\_presentazione.tex).
                    \item Miglioramento del sistema di compilazione automatica dei documenti (compile.sh).
                \end{itemize}


                \paragraph{Rischi attesi}\mbox{}\\
                I componenti del Team ritengono che per questo Sprint siano possibili i seguenti rischi:
                \begin{itemize}
                    \item RC2 - Problemi di comunicazione all’interno del gruppo
                    \item RI2 - Impegni personali e universitari
                \end{itemize}

                \paragraph{Preventivo}\mbox{}\\
                \begin{table}[H]
                    \centering
                    \begin{tabular}{|c|c|c|c|c|c|c|c|}
                    \hline
                                  & Re  & Am  & An  & Pj  & Pg  & Ve  & Totale per persona \\ \hline
                    Cotti Cottini & -   & 3,5 & -   & -   & -   & -   & 3,5                \\ \hline
                    Di Pietro     & -   & -   & -   & -   & -   & 5   & 5                  \\ \hline
                    Diviesti      & -   & -   & -   & -   & 10  & -   & 10                 \\ \hline
                    Lapenna       & -   & -   & -   & -   & -   & 5   & 5                  \\ \hline
                    Pan           & 3   & -   & -   & -   & -   & -   & 3                  \\ \hline
                    Pinarello     & -   & -   & 3,5 & -   & -   & -   & 3,5                \\ \hline
                    Rizzolo       & -   & -   & -   & 5,5 & -   & -   & 5,5                \\ \hline
                    Totale        & 3   & 3,5 & 3,5 & 5,5 & 10  & 10  &                    \\ \hline
                    \end{tabular}
                    \caption{Preventivo orario per ruolo nello Sprint 12}
                \end{table}


                \begin{figure}[H]
                    \centering
                    \begin{tikzpicture}
                        \pie[polar]{8/Responsabile, 10/Amministratore, 10/Analista, 15/Progettista, 28/Programmatore, 28/Verificatore}
                    \end{tikzpicture}
                    \caption{Ripartizione in percentuale dei ruoli nello Sprint 12}
                \end{figure}


                % TODO copiala dal preventivo (per ora li facciamo uguali)
                \paragraph{Consuntivo}\mbox{}\\
                \begin{table}[H]
                    \centering
                    \begin{tabular}{|c|c|c|c|c|c|c|c|}
                    \hline
                                  & Re  & Am  & An  & Pj  & Pg  & Ve  & Totale per persona \\ \hline
                    Cotti Cottini & -   & 3,5 & -   & -   & -   & -   & 3,5                \\ \hline
                    Di Pietro     & -   & -   & -   & -   & -   & 6 \textcolor{red}{(+1)}   & 6 \textcolor{red}{(+1)}                  \\ \hline
                    Diviesti      & -   & -   & -   & -   & 11 \textcolor{red}{(+1)}  & -   & 11 \textcolor{red}{(+1)}                 \\ \hline
                    Lapenna       & -   & -   & -   & -   & -   & 5   & 5                  \\ \hline
                    Pan           & 3   & -   & -   & -   & -   & -   & 3                  \\ \hline
                    Pinarello     & -   & -   & 3,5 & -   & -   & -   & 3,5                \\ \hline
                    Rizzolo       & -   & -   & -   & 6,5 \textcolor{red}{(+1)} & -   & -   & 6,5 \textcolor{red}{(+1)}               \\ \hline
                    Totale        & 3   & 3,5 & 3,5 & 6,5 \textcolor{red}{(+1)} & 11 \textcolor{red}{(+1)}  & 11 \textcolor{red}{(+1)}  &                    \\ \hline
                    \end{tabular}
                    \caption{Consuntivo orario per ruolo nello Sprint 12}
                \end{table}

                % TODO modificare colonna "Ore Rimanenti"
                \paragraph{Prospetto Economico}\mbox{}\\
                \begin{table}[H]
                    \centering
                    \begin{tabular}{|c|c|c|c|}
                    \hline
                    \textbf{Ruolo}  & \textbf{Ore}  & \textbf{Costo} & \textbf{Ore rimanenti} \\ \hline
                    Responsabile    & 3             & €90            & 17                     \\ \hline
                    Amministratore  & 3,5           & €70            & 22                   \\ \hline
                    Analista        & 3,5           & €87,5          & 22                   \\ \hline
                    Progettista     & 6,5           & €162,5         & 37                   \\ \hline
                    Programmatore   & 11            & €165           & 64                     \\ \hline
                    Verificatore    & 11            & €165           & 65                     \\ \hline
                    \textbf{Totale} & \textbf{38,5} & \textbf{\euro740}   & \textbf{227}         \\ \hline
                    \end{tabular}
                    \caption{Prospetto economico e ore rimanenti}
                \end{table}


                \paragraph{Rischi effettivamente occorsi}\mbox{}\\
                Nessuno

                \paragraph{Retrospettiva}
                Risultati Raggiunti:
                Al termine dello sprint, il team ha completato con successo i seguenti elementi:
                \begin{itemize}
                    \item Revisione e aggiornamento dei documenti principali: Piano di Qualifica, Piano di Progetto e Norme di Progetto.
                    \item Correzione e miglioramento dei diagrammi UML (Database e Casi d'Uso).
                    \item Verifica e validazione dei verbali interni e del cliente.
                    \item Implementazione della lettera di presentazione per la documentazione ufficiale.
                    \item Ottimizzazione del sistema di compilazione automatica tramite script.
                \end{itemize}
                
                Miglioramenti per i prossimi Sprint
                Durante lo sprint sono stati riscontrati alcuni punti di miglioramento:
                \begin{itemize}
                    \item Maggiore automazione nei processi di verifica della documentazione.
                    \item Ottimizzazione della gestione dei diagrammi UML per una migliore leggibilità.
                    \item Coinvolgimento più strutturato dei membri del team nella revisione dei documenti.
                \end{itemize}
            %%%%%%%%%%%%%%%%%%%%%%%%%%%%%%%%%%%%%%%%%%%%%%%%%%%%%%%%%%%%%%
            \subsubsection{Sprint 13}
                \paragraph{Pianificazione}
                \begin{itemize}
                    \item inizio: 27 Gennaio 2025
                    \item fine: 2 Febbraio 2025
                \end{itemize}
                Lo sprint si focalizza sull'aggiornamento e il completamento della documentazione del progetto, in particolare su:
                \begin{itemize}
                    \item Aggiornamento dei requisiti del database.
                    \item Revisione e aggiornamento del glossario.
                    \item Aggiornamento dei casi d'uso e del diagramma use-case.
                    \item Verifica e aggiornamento dell'Analisi dei Requisiti.
                    \item Revisione dell'analisi del frontend.
                    \item Verifica e aggiornamento degli ADR (Architecture Decision Records).
                \end{itemize}


                \paragraph{Rischi attesi}\mbox{}\\
                I componenti del Team ritengono che per questo Sprint siano possibili i seguenti rischi:
                \begin{itemize}
                    \item RT1 - Difficoltà nell’uso di nuove tecnologie
                    \item RC2 - Problemi di comunicazione all’interno del gruppo
                    \item RI2 - Impegni personali e universitari
                \end{itemize}

                \paragraph{Preventivo}\mbox{}\\
                \begin{table}[H]
                    \centering
                    \begin{tabular}{|c|c|c|c|c|c|c|c|}
                    \hline
                                  & Re  & Am  & An  & Pj  & Pg  & Ve  & Totale per persona \\ \hline
                    Cotti Cottini & -   & -   & -   & -   & 8  & -   & 8                 \\ \hline
                    Di Pietro     & -   & 3,5 & -   & -   & -   & -   & 3,5                \\ \hline
                    Diviesti      & -   & -   & -   & -   & -   & 5   & 5                  \\ \hline
                    Lapenna       & 3   & -   & -   & -   & -   & -   & 3                  \\ \hline
                    Pan           & -   & -   & -   & 5,5 & -   & -   & 5,5                \\ \hline
                    Pinarello     & -   & -   & -   & -   & -   & 5   & 5                  \\ \hline
                    Rizzolo       & -   & -   & 3,5 & -   & -   & -   & 3,5                \\ \hline
                    Totale        & 3   & 3,5 & 3,5 & 5,5 & 8  & 10  &                    \\ \hline
                    \end{tabular}
                    \caption{Preventivo orario per ruolo nello Sprint 13}
                \end{table}


                \begin{figure}[H]
                    \centering
                    \begin{tikzpicture}
                        \pie[polar]{9/Responsabile, 10/Amministratore, 10/Analista, 16/Progettista, 23/Programmatore, 29/Verificatore}
                    \end{tikzpicture}
                    \caption{Ripartizione in percentuale dei ruoli nello Sprint 13}
                \end{figure}


                % TODO copiala dal preventivo (per ora li facciamo uguali)
                \paragraph{Consuntivo}\mbox{}\\
                \begin{table}[H]
                    \centering
                    \begin{tabular}{|c|c|c|c|c|c|c|c|}
                    \hline
                                  & Re  & Am  & An  & Pj  & Pg  & Ve  & Totale per persona \\ \hline
                    Cotti Cottini & -   & -   & -   & -   & 8  & -   & 8                 \\ \hline
                    Di Pietro     & -   & 3,5 & -   & -   & -   & -   & 3,5                \\ \hline
                    Diviesti      & -   & -   & -   & -   & -   & 6 \textcolor{red}{(+1)}   & 6 \textcolor{red}{(+1)}                  \\ \hline
                    Lapenna       & 3   & -   & -   & -   & -   & -   & 3                  \\ \hline
                    Pan           & -   & -   & -   & 6,5 \textcolor{red}{(+1)} & -   & -   & 6,5 \textcolor{red}{(+1)}               \\ \hline
                    Pinarello     & -   & -   & -   & -   & -   & 7 \textcolor{red}{(+2)}   & 7 \textcolor{red}{(+2)}                  \\ \hline
                    Rizzolo       & -   & -   & 3,5 & -   & -   & -   & 3,5                \\ \hline
                    Totale        & 3   & 3,5 & 3,5 & 6,5 \textcolor{red}{(+1)} & 8  & 13 \textcolor{red}{(+3)}  &                    \\ \hline
                    \end{tabular}
                    \caption{Consuntivo orario per ruolo nello Sprint 13}
                \end{table}

                % TODO modificare colonna "Ore Rimanenti"
                \paragraph{Prospetto Economico}\mbox{}\\
                \begin{table}[H]
                    \centering
                    \begin{tabular}{|c|c|c|c|}
                    \hline
                    \textbf{Ruolo}  & \textbf{Ore}  & \textbf{Costo} & \textbf{Ore rimanenti} \\ \hline
                    Responsabile    & 3             & €90            & 14                     \\ \hline
                    Amministratore  & 3,5           & €70            & 18,5                     \\ \hline
                    Analista        & 3,5           & €87,5          & 18,5                     \\ \hline
                    Progettista     & 6,5           & €162,5         & 30,5                     \\ \hline
                    Programmatore   & 8            & €120           & 56                     \\ \hline
                    Verificatore    & 13            & €195           & 52                     \\ \hline
                    \textbf{Totale} & \textbf{37,5} & \textbf{\euro755}   & \textbf{189,5}         \\ \hline
                    \end{tabular}
                    \caption{Prospetto economico e ore rimanenti}
                \end{table}


                \paragraph{Rischi effettivamente occorsi}\mbox{}\\
                RC2 - Problemi di comunicazione all’interno del gruppo \\
                Il problema principale causato da questo rischio riguarda \textit{commit} ridondanti e sovrapposizioni nelle modifiche che si possono verificare nei documenti. \\
                \underline{Possibili migliorie:} Oltre ad aver stabilito i canali specifici di comunicazione e alle riunioni svolte settimanalmente, è importante definire nello specifico le responsabilità dei membri del Team e ottimizzare il workflow di aggiornamento dei vari documenti

                \paragraph{Retrospettiva}
                Durante lo sprint, il team ha completato con successo le seguenti attività:
                \begin{itemize}
                    \item Eliminazione di PDF inutili e aggiornamento dei requisiti del database.
                    \item Revisione e aggiornamento del glossario.
                    \item Aggiornamento dei casi d'uso sulla base delle indicazioni post revisione.
                    \item Modifica del diagramma use-case in base agli ultimi aggiornamenti.
                    \item Revisione dell'analisi frontend per una migliore strutturazione della documentazione.
                    \item Fix di problemi minori nei percorsi dei file.
                    \item Verifica e approvazione dell'ADR numero 117.
                \end{itemize}
                
                Feedback e Miglioramenti per i Prossimi Sprint:
                \begin{itemize}
                    \item Ottimizzare il workflow di aggiornamento della documentazione per evitare commit ridondanti.
                    \item Definire in modo più chiaro le responsabilità per evitare sovrapposizioni nelle modifiche.
                    \item Automatizzare ulteriormente il processo di verifica dei documenti prima della generazione del PDF.
                \end{itemize}
            %%%%%%%%%%%%%%%%%%%%%%%%%%%%%%%%%%%%%%%%%%%%%%%%%%%%%%%%%%%%%%
            \subsubsection{Sprint 14}
                \paragraph{Pianificazione}
                \begin{itemize}
                    \item inizio: 3 Febbraio 2025
                    \item fine: 9 Febbraio 2025
                \end{itemize}
                Durante questo sprint, il team si è prefissato di completare le seguenti attività:
                \begin{itemize}
                    \item Revisione e aggiornamento dei documenti di progetto, tra cui:
                    \begin{itemize}
                        \item Revisione e aggiornamento del Piano di Qualifica (\texttt{piano\_di\_qualifica\_v1.0.tex}).
                        \item Revisione e aggiornamento del Piano di Progetto (\texttt{piano\_di\_progetto\_v1.0.tex}).
                        \item Correzione delle Norme di Progetto.
                    \end{itemize}
                    \item Verifica e aggiornamento della documentazione:
                    \begin{itemize}
                        \item Verifica e aggiornamento del Glossario.
                        \item Verifica dei verbali interni ed esterni:
                        \begin{itemize}
                            \item Verbale interno del 24 gennaio 2025.
                            \item Verbale esterno del 27 gennaio 2025.
                            \item Verbale interno del 18 gennaio 2025.
                            \item Verbale interno del 7 febbraio 2025.
                        \end{itemize}
                    \end{itemize}
                    \item Correzioni grafiche e tecniche per migliorare la qualità dei documenti generati (fix di grafica e generazione PDF).
                    \item Merge delle modifiche nel branch di sviluppo.
                \end{itemize}


                \paragraph{Rischi attesi}\mbox{}\\
                I componenti del Team ritengono che per questo Sprint siano possibili i seguenti rischi:
                \begin{itemize}
                    \item RC3 - Mancato confronto con l’azienda Proponente
                    \item RO2 - Deviazione dai tempi e costi previsti
                    \item RO1 - Incomprensioni nella pianificazione delle attività
                \end{itemize}

                \paragraph{Preventivo}\mbox{}\\
                \begin{table}[H]
                    \centering
                    \begin{tabular}{|c|c|c|c|c|c|c|c|}
                    \hline
                                  & Re  & Am  & An  & Pj  & Pg  & Ve  & Totale per persona \\ \hline
                    Cotti Cottini & -   & -   & -   & -   & -   & 5   & 5                  \\ \hline
                    Di Pietro     & -   & -   & -   & -   & 10  & -   & 10                 \\ \hline
                    Diviesti      & 3   & -   & -   & -   & -   & -   & 3                  \\ \hline
                    Lapenna       & -   & -   & -   & 5,5 & -   & -   & 5,5                \\ \hline
                    Pan           & -   & -   & 3,5 & -   & -   & -   & 3,5                \\ \hline
                    Pinarello     & -   & 3,5 & -   & -   & -   & -   & 3,5                \\ \hline
                    Rizzolo       & -   & -   & -   & -   & -   & 5   & 5                  \\ \hline
                    Totale        & 3   & 3,5 & 3,5 & 5,5 & 10  & 10  &                    \\ \hline
                    \end{tabular}
                    \caption{Preventivo orario per ruolo nello Sprint 14}
                \end{table}


                \begin{figure}[H]
                    \centering
                    \begin{tikzpicture}
                        \pie[polar]{8/Responsabile, 10/Amministratore, 10/Analista, 16/Progettista, 28/Programmatore, 28/Verificatore}
                    \end{tikzpicture}
                    \caption{Ripartizione in percentuale dei ruoli nello Sprint 14}
                \end{figure}


                % TODO copiala dal preventivo (per ora li facciamo uguali)
                \paragraph{Consuntivo}\mbox{}\\
                \begin{table}[H]
                    \centering
                    \begin{tabular}{|c|c|c|c|c|c|c|c|}
                    \hline
                                  & Re  & Am  & An  & Pj  & Pg  & Ve  & Totale per persona \\ \hline
                    Cotti Cottini & -   & -   & -   & -   & -   & 6 \textcolor{red}{(+1)}   & 6 \textcolor{red}{(+1)}                  \\ \hline
                    Di Pietro     & -   & -   & -   & -   & 10  & -   & 10                 \\ \hline
                    Diviesti      & 3   & -   & -   & -   & -   & -   & 3                  \\ \hline
                    Lapenna       & -   & -   & -   & 6,5 \textcolor{red}{(+1)} & -   & -   & 6,5 \textcolor{red}{(+1)}               \\ \hline
                    Pan           & -   & -   & 4,5 \textcolor{red}{(+1)} & -   & -   & -   & 4,5 \textcolor{red}{(+1)}               \\ \hline
                    Pinarello     & -   & 3,5 & -   & -   & -   & -   & 3,5                \\ \hline
                    Rizzolo       & -   & -   & -   & -   & -   & 7 \textcolor{red}{(+2)}   & 7 \textcolor{red}{(+2)}                  \\ \hline
                    Totale        & 3   & 3,5 & 4,5 \textcolor{red}{(+1)} & 6,5 \textcolor{red}{(+1)} & 10  & 13 \textcolor{red}{(+3)}  &                    \\ \hline
                    \end{tabular}
                    \caption{Consuntivo orario per ruolo nello Sprint 14}
                \end{table}

                % TODO modificare colonna "Ore Rimanenti"
                \paragraph{Prospetto Economico}\mbox{}\\
                \begin{table}[H]
                    \centering
                    \begin{tabular}{|c|c|c|c|}
                    \hline
                    \textbf{Ruolo}  & \textbf{Ore}  & \textbf{Costo} & \textbf{Ore rimanenti} \\ \hline
                    Responsabile    & 3             & €90            & 11                     \\ \hline
                    Amministratore  & 3,5           & €70            & 15                   \\ \hline
                    Analista        & 4,5           & €112,5          & 14                   \\ \hline
                    Progettista     & 6,5           & €162,5         & 24                   \\ \hline
                    Programmatore   & 10            & €150           & 46                     \\ \hline
                    Verificatore    & 13            & €195           & 39                     \\ \hline
                    \textbf{Totale} & \textbf{40,5} & \textbf{\euro780}   & \textbf{149}         \\ \hline
                    \end{tabular}
                    \caption{Prospetto economico e ore rimanenti}
                \end{table}


                \paragraph{Rischi effettivamente occorsi}\mbox{}\\
                Nessuno

                \paragraph{Retrospettiva}
                Al termine dello sprint, il team ha raggiunto i seguenti risultati:
                \begin{itemize}
                    \item Sono stati aggiornati e verificati i principali documenti di progetto.
                    \item I verbali sono stati verificati, rinominati e corretti.
                    \item Sono stati risolti problemi tecnici legati alla generazione dei PDF.
                    \item È stato effettuato il merge delle modifiche nel branch di sviluppo.
                \end{itemize}
                
                Dopo la revisione dello sprint, il team ha identificato alcune aree di miglioramento:
                \begin{itemize}
                    \item Ottimizzazione del processo di verifica dei documenti per evitare sovrapposizioni nelle revisioni.
                    \item Automazione avanzata per il controllo di qualità dei documenti prima della generazione dei PDF.
                    \item Maggiore attenzione alla gestione delle versioni dei documenti per evitare incongruenze.
                \end{itemize}
            %%%%%%%%%%%%%%%%%%%%%%%%%%%%%%%%%%%%%%%%%%%%%%%%%%%%%%%%%%%%%%
            \subsubsection{Sprint 15}
                \paragraph{Pianificazione}
                \begin{itemize}
                    \item inizio: 10 Febbraio 2025
                    \item fine: 16 Febbraio 2025
                \end{itemize}

                \paragraph{Obiettivi dello Sprint}
                    \begin{itemize}
                        \item Approvazione Norme di Progetto
                        \item Motivazioni delle tecnologie utilizzate.
                        \item Sistemazione Use Case in Analisi dei Requisiti.
                    \end{itemize}

                \paragraph{Task Prioritari}
                    \begin{itemize}
                        \item Motivare le tecnologie usate nel PoC.
                    \end{itemize}                


                \paragraph{Rischi attesi}\mbox{}\\
                    I componenti del Team ritengono che per questo Sprint siano possibili i seguenti rischi:
                    \begin{itemize}
                        \item RT1 - Difficoltà nell’uso di nuove tecnologie
                        \item RI2 - Impegni personali e universitari
                    \end{itemize}


                \paragraph{Preventivo}\mbox{}\\
                \begin{table}[H]
                    \centering
                    \begin{tabular}{|c|c|c|c|c|c|c|c|}
                    \hline
                                  & Re  & Am  & An  & Pj  & Pg  & Ve  & Totale per persona \\ \hline
                    Cotti Cottini & 2   & -   & -   & -   & -   & -   & 2                  \\ \hline
                    Di Pietro     & -   & -   & -   & -   & -   & 5   & 5                  \\ \hline
                    Diviesti      & -   & -   & -   & 5,5 & -   & -   & 5,5                \\ \hline
                    Lapenna       & -   & -   & 3,5 & -   & -   & -   & 3,5                \\ \hline
                    Pan           & -   & -   & -   & -   & -   & 5   & 5                  \\ \hline
                    Pinarello     & -   & -   & -   & -   & 8  & -   & 8                 \\ \hline
                    Rizzolo       & -   & 3,5 & -   & -   & -   & -   & 3,5                \\ \hline
                    Totale        & 2   & 3,5 & 3,5 & 5,5 & 8  & 10  &                    \\ \hline
                    \end{tabular}
                    \caption{Preventivo orario per ruolo nello Sprint 15}
                \end{table}


                \begin{figure}[H]
                    \centering
                    \begin{tikzpicture}
                        \pie[polar]{6/Responsabile, 11/Amministratore, 11/Analista, 17/Progettista, 25/Programmatore, 31/Verificatore}
                    \end{tikzpicture}
                    \caption{Ripartizione in percentuale dei ruoli nello Sprint 15}
                \end{figure}


                % TODO copiala dal preventivo (per ora li facciamo uguali)
                \paragraph{Consuntivo}\mbox{}\\
                \begin{table}[H]
                    \centering
                    \begin{tabular}{|c|c|c|c|c|c|c|c|}
                    \hline
                                  & Re  & Am  & An  & Pj  & Pg  & Ve  & Totale per persona \\ \hline
                    Cotti Cottini & 2   & -   & -   & -   & -   & -   & 2                  \\ \hline
                    Di Pietro     & -   & -   & -   & -   & -   & 7 \textcolor{red}{(+2)}   & 7 \textcolor{red}{(+2)}                  \\ \hline
                    Diviesti      & -   & -   & -   & 7,5 \textcolor{red}{(+2)} & -   & -   & 7,5 \textcolor{red}{(+2)}                \\ \hline
                    Lapenna       & -   & -   & 3,5 & -   & -   & -   & 3,5                \\ \hline
                    Pan           & -   & -   & -   & -   & -   & 6 \textcolor{red}{(+1)}   & 6 \textcolor{red}{(+1)}                  \\ \hline
                    Pinarello     & -   & -   & -   & -   & 8  & -   & 8                 \\ \hline
                    Rizzolo       & -   & 3,5 & -   & -   & -   & -   & 3,5                \\ \hline
                    Totale        & 2   & 3,5 & 3,5 & 7,5 \textcolor{red}{(+2)} & 8  & 13 \textcolor{red}{(+3)}  &                    \\ \hline
                    \end{tabular}
                    \caption{Preventivo orario per ruolo nello Sprint 15}  
                \end{table}

                % TODO modificare colonna "Ore Rimanenti"
                \paragraph{Prospetto Economico}\mbox{}\\
                \begin{table}[H]
                    \centering
                    \begin{tabular}{|c|c|c|c|}
                    \hline
                    \textbf{Ruolo}  & \textbf{Ore}  & \textbf{Costo} & \textbf{Ore rimanenti} \\ \hline
                    Responsabile    & 2             & €60            & 9                     \\ \hline
                    Amministratore  & 3,5           & €70            & 11,5                     \\ \hline
                    Analista        & 3,5           & €87,5          & 10,5                     \\ \hline
                    Progettista     & 7,5           & €187,5         & 16,5                     \\ \hline
                    Programmatore   & 8            & €120           & 38                     \\ \hline
                    Verificatore    & 13            & €195           & 26                     \\ \hline
                    \textbf{Totale} & \textbf{37,5} & \textbf{\euro735}   & \textbf{111,5}           \\ \hline
                    \end{tabular}
                    \caption{Prospetto economico e ore rimanenti}
                \end{table}


                \paragraph{Rischi effettivamente occorsi}\mbox{}\\
                RT1 - Difficoltà nell’uso di nuove tecnologie \\
                Nonostante lo studio precedente e la stesura di varie documentazioni a riguardo, è stato utilizzato del tempo per comprendere al meglio delle scelte tecnologiche e approfondire diverse questioni che le riguardano, come per esempio l’integrazione di un database vettoriale. \\
                Come già supposto durante l’analisi dei rischi, dato che si tratta di tecnologie che non conosciamo alla perfezione, è normale e buona norma impiegare del tempo per capire bene il loro funzionamento. \\
                \underline{Possibili migliorie:} Nessuna

                \paragraph{Retrospettiva}
                Durante lo sprint sono stati completati con successo i seguenti obiettivi:
                \begin{itemize}
                    \item \textbf{Caso d'Uso U.C.20:} Il caso d'uso U.C.20 è stato verificato e i risultati sono stati inclusi nel documento di analisi.
                    \item \textbf{Documentazione dei Requisiti:} I requisiti per il chatbot e la gestione delle metriche sono stati aggiornati nel file \texttt{requisiti.tex}.
                    \item \textbf{Analisi dei Database Vettoriali:} È stata completata una nuova sezione introduttiva sui database vettoriali, con particolare attenzione a FAISS e un confronto con altri strumenti come Annoy e Milvus.
                    \item \textbf{Compilazione Automatica:} Sono stati apportati miglioramenti allo script di compilazione \texttt{compile.sh}, con l'inclusione di file precedentemente commentati.
                    \item \textbf{Conclusioni su LangChain:} Sono state scritte conclusioni sui vantaggi e svantaggi dell'uso di LangChain per la gestione dei flussi conversazionali avanzati.
                \end{itemize}
                
                Non sono stati riscontrati problemi critici durante lo sprint, e tutte le attività sono state completate secondo il piano. Tuttavia, l'integrazione con alcuni database vettoriali ha richiesto un approfondimento delle loro caratteristiche per garantire la miglior scelta in base ai requisiti di progetto.
                
                Prossimi Passi:
                \begin{itemize}
                    \item Continuare l'analisi e l'implementazione del chatbot, con particolare attenzione all'integrazione con modelli di embedding.
                    \item Espandere la documentazione relativa alla gestione del database, includendo dettagli sulle prestazioni e la scalabilità.
                    \item Continuare l'analisi delle librerie per flussi conversazionali avanzati, esplorando alternative a LangChain.
                \end{itemize}


                \subsubsection{Retrospettiva di Periodo (RTB)}

                Allo stato attuale (fine dello Sprint 15) sono ancora disponibili 
                \textbf{111,5 ore} per il completamento del progetto. Fino a questo punto, 
                sono state contabilizzate un totale di \textbf{541,5 ore} lavorative, 
                suddivise tra i diversi ruoli del team. Il consuntivo economico 
                registrato fino ad ora ammonta a \textbf{10.530 euro}. I costi e i tempi 
                sostenuti finora risultano in linea con le previsioni iniziali, 
                senza alcun significativo scostamento rispetto a quanto pianificato. 
                Questo indica che il progetto sta procedendo secondo le tempistiche 
                e i budget stabiliti, con una gestione efficace delle risorse e un 
                buon allineamento con gli obiettivi prefissati. In sintesi, i dati 
                attuali confermano che il progetto è sotto controllo e sta 
                rispettando le aspettative sia dal punto di vista economico che 
                temporale.


                \definecolor{re}{HTML}{6666f6}  
                \definecolor{am}{HTML}{77d2f3}  
                \definecolor{an}{HTML}{fdf399}  
                \definecolor{pj}{HTML}{f4b572}  
                \definecolor{pg}{HTML}{ed6f6b}  
                \definecolor{ve}{HTML}{74a0f8} 

                \begin{figure}[H]
                    \centering
                    \begin{tikzpicture}
                        \pie[polar]{8.3/Responsabile, 9.7/Amministratore, 10.2/Analista, 16.9/Progettista, 22.2/Programmatore, 32.7/Verificatore}
                    \end{tikzpicture}
                    \caption{Lavoro in percentuale dei ruoli nel periodo RTB}
                \end{figure}

                \vspace{1cm}

                \begin{figure}[H]
                    \centering
                    \begin{tikzpicture}
                        \begin{axis}[
                            width=1.1*\textwidth,
                            height=8cm,
                            ybar,
                            bar width=20pt,
                            ymajorgrids=true,
                            symbolic x coords={p1, p2, p3, p4, p5, p6},
                            scaled y ticks=false,
                            enlarge x limits=0.3,
                            ymin=0,
                            enlarge y limits=0.1,
                            legend cell align=left,
                            legend style={
                                inner sep=8pt, 
                                at={(0.5,1.35)},
                                anchor=north,
                                column sep=8pt,
                                row sep=5pt,
                                legend columns=3,
                                font=\large,
                            },
                            xticklabel=\empty,
                            ylabel={Euro (\euro)}
                        ] 
                        % Responsabile (1350) 
                        \addplot[fill=re] coordinates {(p1,1350) (p2,0) (p3,0) (p4,0) (p5,0) (p6,0)};  
                        % Amministratore (1050) 
                        \addplot[fill=am] coordinates {(p1,0) (p2,1050) (p3,0) (p4,0) (p5,0) (p6,0)};  
                        % Analista (1388) 
                        \addplot[fill=an] coordinates {(p1,0) (p2,0) (p3,1388) (p4,0) (p5,0) (p6,0)};  
                        % Progettista (2288) 
                        \addplot[fill=pj] coordinates {(p1,0) (p2,0) (p3,0) (p4,2288) (p5,0) (p6,0)};  
                        % Programmatore (1800)
                        \addplot[fill=pg] coordinates {(p1,0) (p2,0) (p3,0) (p4,0) (p5,1800) (p6,0)};  
                        % Verificatore (2655)
                        \addplot[fill=ve] coordinates {(p1,0) (p2,0) (p3,0) (p4,0) (p5,0) (p6,2655)};  
                
                        \legend{Responsabile, Amministratore, Analista, Progettista, Programmatore, Verificatore}
                        \end{axis}
                    \end{tikzpicture}
                    \caption{Consuntivo economico per ruolo nel periodo RTB}
                \end{figure}



            %%%%%%%%%%%%%%%%%%%%%%%%%%%%%%%%%%%%%%%%%%%%%%%%%%%%%%%%%%%%%%
        \subsection{PB}
            %%%%%%%%%%%%%%%%%%%%%%%%%%%%%%%%%%%%%%%%%%%%%%%%%%%%%%%%%%%%%%
            \subsubsection{Sprint 16}
            \paragraph{Pianificazione}
                \begin{itemize}
                    \item inizio: 24 Febbraio 2025
                    \item fine: 2 Marzo 2025
                \end{itemize}
            
            \paragraph{Obbiettivi dello Sprint}
            Durante questo sprint, il team ha lavorato principalmente sulla creazione del verbale interno, il glossario, il manuale e la specifica, nonché sull'aggiornamento del tracciamento dei casi d’uso e dei requisiti funzionali. Inoltre, è stato aggiunto il caso d’uso U.C.23 relativo alla visualizzazione del Template di Domanda e Risposta.
            
            \paragraph{Task Prioritari}
            \begin{itemize}
                \item Creazione del verbale interno del 28 febbraio
                \item Inizializzazione del Manuale e della Specifica
                \item Aggiornamento del glossario
                \item Versionamento della documentazione alla v2.1
                \item Tracciamento aggiornato dei requisiti funzionali
                \item Aggiunta del caso d’uso U.C.23
            \end{itemize}
            
            \paragraph{Rischi attesi}
            Sulla base delle attività pianificate, i rischi principali ipotizzati per lo sprint erano:
            \begin{itemize}
                \item RT1: Difficoltà nell’uso di nuove tecnologie
                \item RT2: Sviluppo di codice poco comprensibile
                \item RC2: Problemi di comunicazione all’interno del gruppo
                \item RO2: Deviazione dai tempi e costi previsti
            \end{itemize}

            \paragraph{Preventivo}\mbox{}\\
            \begin{table}[H]
                \centering
                \begin{tabular}{|c|c|c|c|c|c|c|c|}
                \hline
                              & Re  & Am  & An  & Pj  & Pg  & Ve  & Totale per persona \\ \hline
                Cotti Cottini & -   & -   & -   & 2,5 & -   & -   & 2,5                \\ \hline
                Di Pietro     & 1   & -   & -   & -   & -   & -   & 1                  \\ \hline
                Diviesti      & -   & -   & 2   & -   & -   & -   & 2                  \\ \hline
                Lapenna       & -   & -   & -   & -   & -   & 3   & 3                  \\ \hline
                Pan           & -   & 2   & -   & -   & -   & -   & 2                  \\ \hline
                Pinarello     & -   & -   & -   & -   & -   & 3   & 3                  \\ \hline
                Rizzolo       & -   & -   & -   & -   & 6,5 & -   & 6,5                \\ \hline
                Totale        & 1   & 2   & 2   & 2,5 & 6,5 & 6   &                    \\ \hline
                \end{tabular}
                \caption{Preventivo orario per ruolo nello Sprint 16}
            \end{table}


            \begin{figure}[H]
                \centering
                \begin{tikzpicture}
                    \pie[polar]{5/Responsabile, 10/Amministratore, 10/Analista, 12/Progettista, 33/Programmatore, 30/Verificatore}
                \end{tikzpicture}
                \caption{Ripartizione in percentuale dei ruoli nello Sprint 16}
            \end{figure}


            % TODO copiala dal preventivo (per ora li facciamo uguali)
            \paragraph{Consuntivo}\mbox{}\\
            \begin{table}[H]
                \centering
                \begin{tabular}{|c|c|c|c|c|c|c|c|}
                \hline
                              & Re  & Am  & An  & Pj  & Pg  & Ve  & Totale per persona \\ \hline
                Cotti Cottini & -   & -   & -   & 2,5 & -   & -   & 2,5                \\ \hline
                Di Pietro     & 1   & -   & -   & -   & -   & -   & 1                  \\ \hline
                Diviesti      & -   & -   & 2   & -   & -   & -   & 2                  \\ \hline
                Lapenna       & -   & -   & -   & -   & -   & 3   & 3                  \\ \hline
                Pan           & -   & 2   & -   & -   & -   & -   & 2                  \\ \hline
                Pinarello     & -   & -   & -   & -   & -   & 3   & 3                  \\ \hline
                Rizzolo       & -   & -   & -   & -   & 6,5 & -   & 6,5                \\ \hline
                Totale        & 1   & 2   & 2   & 2,5 & 6,5 & 6   &                    \\ \hline
                \end{tabular}
                \caption{Consuntivo orario per ruolo nello Sprint 16}
            \end{table}

            % TODO modificare colonna "Ore Rimanenti"
            \paragraph{Prospetto Economico}\mbox{}\\
            \begin{table}[H]
                \centering
                \begin{tabular}{|c|c|c|c|}
                \hline
                \textbf{Ruolo}  & \textbf{Ore}  & \textbf{Costo}      & \textbf{Ore rimanenti}   \\ \hline
                Responsabile    & 1             & €30                 & 8                        \\ \hline
                Amministratore  & 2             & €40                 & 9,5                      \\ \hline
                Analista        & 2             & €50                 & 8,5                      \\ \hline
                Progettista     & 2,5           & €62,5               & 14                       \\ \hline
                Programmatore   & 6,5           & €97,5               & 31,5                     \\ \hline
                Verificatore    & 6             & €90                 & 20                       \\ \hline
                \textbf{Totale} & \textbf{20}   & \textbf{\euro370}   & \textbf{91,5}            \\ \hline
                \end{tabular}
                \caption{Prospetto economico e ore rimanenti}
            \end{table}
            
            \paragraph{Rischi effettivamente occorsi}
            Durante lo sprint, sono stati riscontrati i seguenti rischi:
            \begin{itemize}
                \item RT1: Difficoltà nell’uso di nuove tecnologie, in particolare nell'aggiornamento del pacchetto \texttt{pgfplots} e nella gestione della documentazione.
                \item RO2: Alcune attività hanno subito ritardi rispetto alla pianificazione iniziale, in particolare la creazione della specifica tecnica.
            \end{itemize}
            
            \paragraph{Retrospettiva}
            Nel complesso, lo sprint ha portato a risultati significativi, tra cui l'inizializzazione del manuale e della specifica, l'aggiornamento della documentazione e l'aggiunta di nuovi casi d'uso. Tuttavia, sono emersi alcune difficoltà nell'adattarsi a nuovi strumenti, che in parte hanno rallentato i progressi. Per migliorare l'efficienza nei prossimi sprint, si suggerisce di:
            \begin{itemize}
                \item Definire chiaramente le responsabilità di ogni membro del team per ridurre le incertezze.
                \item Effettuare una formazione più approfondita sugli strumenti e le tecnologie utilizzate nel progetto.
            \end{itemize}
                
                
            %%%%%%%%%%%%%%%%%%%%%%%%%%%%%%%%%%%%%%%%%%%%%%%%%%%%%%%%%%%%%%
            \subsubsection{Sprint 17}
            \paragraph{Pianificazione}
                \begin{itemize}
                    \item inizio: 3 Marzo 2025
                    \item fine: 9 Marzo 2025
                \end{itemize}
            
            \paragraph{Obbiettivi dello Sprint}
            In questo sprint, il team si è concentrato principalmente sulla gestione e revisione della documentazione, con particolare attenzione ai verbali interni ed esterni, al glossario e agli aggiornamenti relativi al database e ai diagrammi UML.
            
            \paragraph{Task Prioritari}
            \begin{itemize}
                \item Creazione del verbale interno del 4 marzo 2025
                \item Revisione e sistemazione degli indici nei verbali
                \item Eliminazione di file PDF obsoleti
                \item Aggiornamento del glossario
                \item Aggiunta della sezione 4.9 AdR nella documentazione
                \item Integrazione di dettagli su chat nei documenti
                \item Aggiunta di un diagramma UML esagonale
                \item Inserimento di un file StarUML per il diagramma delle classi
                \item Aggiornamento del file \texttt{db.sql}
            \end{itemize}
            
            \paragraph{Rischi attesi}
            \begin{itemize}
                \item RT1: Difficoltà nell’uso di nuove tecnologie (gestione di diagrammi UML e aggiornamenti del database)
                \item RO1: Incomprensioni nella pianificazione delle attività
            \end{itemize}

            \paragraph{Preventivo}\mbox{}\\
            \begin{table}[H]
                \centering
                \begin{tabular}{|c|c|c|c|c|c|c|c|}
                \hline
                              & Re  & Am  & An  & Pj  & Pg  & Ve  & Totale per persona \\ \hline
                Cotti Cottini & -   & -   & 1,5 & -   & -   & -   & 1,5                \\ \hline
                Di Pietro     & -   & -   & -   & 2,5 & -   & -   & 2,5                \\ \hline
                Diviesti      & -   & -   & -   & -   & -   & 3   & 3                  \\ \hline
                Lapenna       & -   & 1,5 & -   & -   & -   & -   & 1,5                \\ \hline
                Pan           & -   & -   & -   & -   & 6   & -   & 6                  \\ \hline
                Pinarello     & 1   & -   & -   & -   & -   & -   & 1                  \\ \hline
                Rizzolo       & -   & -   & -   & -   & -   & 3   & 3                  \\ \hline
                Totale        & 1   & 1,5 & 1,5 & 2,5 & 6   & 6   &                    \\ \hline
                \end{tabular}
                \caption{Preventivo orario per ruolo nello Sprint 17}
            \end{table}


            \begin{figure}[H]
                \centering
                \begin{tikzpicture}
                    \pie[polar]{5/Responsabile, 8/Amministratore, 8/Analista, 14/Progettista, 33/Programmatore, 32/Verificatore}
                \end{tikzpicture}
                \caption{Ripartizione in percentuale dei ruoli nello Sprint 17}
            \end{figure}


            % TODO copiala dal preventivo (per ora li facciamo uguali)
            \paragraph{Consuntivo}\mbox{}\\
            \begin{table}[H]
                \centering
                \begin{tabular}{|c|c|c|c|c|c|c|c|}
                \hline
                              & Re  & Am  & An  & Pj  & Pg  & Ve  & Totale per persona \\ \hline
                Cotti Cottini & -   & -   & 1,5 & -   & -   & -   & 1,5                \\ \hline
                Di Pietro     & -   & -   & -   & 2,5 & -   & -   & 2,5                \\ \hline
                Diviesti      & -   & -   & -   & -   & -   & 3   & 3                  \\ \hline
                Lapenna       & -   & 1,5 & -   & -   & -   & -   & 1,5                \\ \hline
                Pan           & -   & -   & -   & -   & 6   & -   & 6                  \\ \hline
                Pinarello     & 1   & -   & -   & -   & -   & -   & 1                  \\ \hline
                Rizzolo       & -   & -   & -   & -   & -   & 3   & 3                  \\ \hline
                Totale        & 1   & 1,5 & 1,5 & 2,5 & 6   & 6   &                    \\ \hline
                \end{tabular}
                \caption{Consuntivo orario per ruolo nello Sprint 17}
            \end{table}

            % TODO modificare colonna "Ore Rimanenti"
            \paragraph{Prospetto Economico}\mbox{}\\
            \begin{table}[H]
                \centering
                \begin{tabular}{|c|c|c|c|}
                \hline
                \textbf{Ruolo}  & \textbf{Ore}  & \textbf{Costo}      & \textbf{Ore rimanenti}   \\ \hline
                Responsabile    & 1             & €30                 & 7                        \\ \hline
                Amministratore  & 1,5           & €30                 & 8                        \\ \hline
                Analista        & 1,5           & €37,5               & 7                        \\ \hline
                Progettista     & 2,5           & €62,5               & 11,5                     \\ \hline
                Programmatore   & 6             & €90                 & 25,5                     \\ \hline
                Verificatore    & 6             & €90                 & 14                       \\ \hline
                \textbf{Totale} & \textbf{20}   & \textbf{\euro340}   & \textbf{73}              \\ \hline
                \end{tabular}
                \caption{Prospetto economico e ore rimanenti}
            \end{table}
            
            \paragraph{Rischi effettivamente occorsi}
            \begin{itemize}
                \item RT1: Alcune difficoltà nell'uso di strumenti per la modellazione UML e la gestione del database, che hanno richiesto più tempo del previsto per essere integrate nella documentazione.
                \item RO1: Incomprensioni nella gestione dei verbali hanno reso necessarie correzioni multiple sugli indici e sulla struttura dei documenti.
            \end{itemize}
            
            \paragraph{Retrospettiva}
            Lo sprint ha portato a un miglioramento significativo nella documentazione e nell'organizzazione dei file, con la creazione e revisione di diversi documenti chiave. Tuttavia, si sono verificati problemi di comunicazione e difficoltà tecniche che hanno richiesto più iterazioni per essere risolti. Per migliorare nei prossimi sprint, si consiglia di:
            \begin{itemize}
                \item Coordinarsi meglio per evitare revisioni ripetitive della documentazione.
                \item Definire una strategia più chiara per la gestione degli strumenti di modellazione UML e del database.
                \item Allocare più tempo per la fase di pianificazione, per ridurre gli errori nella gestione dei verbali e dei documenti.
            \end{itemize}
                
            %%%%%%%%%%%%%%%%%%%%%%%%%%%%%%%%%%%%%%%%%%%%%%%%%%%%%%%%%%%%%%
            \subsubsection{Sprint 18}
            \paragraph{Pianificazione}
                \begin{itemize}
                    \item inizio: 10 Marzo 2025
                    \item fine: 16 Marzo 2025
                \end{itemize}
        
            \paragraph{Obbiettivi dello Sprint}
            Durante questo sprint, il team si è focalizzato su diversi aspetti del progetto, tra cui:
            \begin{itemize}
                \item Miglioramento e verifica della documentazione tecnica, inclusi il Piano di Qualifica e il Piano di Progetto.
                \item Integrazione di nuove funzionalità nel backend, come il repository per FAISS e la gestione della memoria utente nella conversazione.
                \item Avanzamenti nella gestione dei database, con l'implementazione del repository PostgreSQL e test di connessione.
                \item Sviluppo del frontend, con la configurazione iniziale e l’integrazione del routing.
            \end{itemize}
            
            \paragraph{Task Prioritari}
            \begin{itemize}
                \item Creazione del diario di bordo del 17 marzo 2025.
                \item Aggiornamenti al Piano di Progetto e al Piano di Qualifica.
                \item Verifica e correzione della Specificha Tecnica (v0.1 e v0.2).
                \item Integrazione e miglioramento del repository FAISS per la similarità.
                \item Implementazione della gestione dei file e della suddivisione dei chunk di dati.
                \item Configurazione dell’ambiente di test con Pytest.
                \item Sviluppo della logica di risposta nel repository LangChain.
                \item Configurazione e test di connessione al database PostgreSQL.
                \item Impostazione del routing e della struttura iniziale per il frontend.
            \end{itemize}
            
            \paragraph{Rischi attesi}
            \begin{itemize}
                \item RT1: Difficoltà nell’uso di nuove tecnologie (FAISS, LangChain, PostgreSQL).
                \item RT2: Sviluppo di codice poco comprensibile, soprattutto nelle parti più complesse come il repository LangChain.
                \item RO1: Incomprensioni nella pianificazione delle attività, con possibili ritardi nell’integrazione delle varie componenti.
            \end{itemize}

            \paragraph{Preventivo}\mbox{}\\
            \begin{table}[H]
                \centering
                \begin{tabular}{|c|c|c|c|c|c|c|c|}
                \hline
                              & Re  & Am  & An  & Pj  & Pg  & Ve  & Totale per persona \\ \hline
                Cotti Cottini & -   & -   & -   & -   & -   & 3   & 3                  \\ \hline
                Di Pietro     & -   & -   & 2   & -   & -   & -   & 2                  \\ \hline
                Diviesti      & -   & 2   & -   & -   & -   & -   & 2                  \\ \hline
                Lapenna       & -   & -   & -   & -   & 6   & -   & 6                  \\ \hline
                Pan           & -   & -   & -   & -   & -   & 3   & 3                  \\ \hline
                Pinarello     & -   & -   & -   & 3   & -   & -   & 3                  \\ \hline
                Rizzolo       & 1,5 & -   & -   & -   & -   & -   & 1,5                \\ \hline
                Totale        & 1,5 & 2   & 2   & 3   & 6   & 6   &                    \\ \hline
                \end{tabular}
                \caption{Preventivo orario per ruolo nello Sprint 18}
            \end{table}


            \begin{figure}[H]
                \centering
                \begin{tikzpicture}
                    \pie[polar]{7/Responsabile, 10/Amministratore, 10/Analista, 14/Progettista, 29/Programmatore, 30/Verificatore}
                \end{tikzpicture}
                \caption{Ripartizione in percentuale dei ruoli nello Sprint 18}
            \end{figure}


            % TODO copiala dal preventivo (per ora li facciamo uguali)
            \paragraph{Consuntivo}\mbox{}\\
            \begin{table}[H]
                \centering
                \begin{tabular}{|c|c|c|c|c|c|c|c|}
                \hline
                              & Re  & Am  & An  & Pj  & Pg  & Ve  & Totale per persona \\ \hline
                Cotti Cottini & -   & -   & -   & -   & -   & 3   & 3                  \\ \hline
                Di Pietro     & -   & -   & 2   & -   & -   & -   & 2                  \\ \hline
                Diviesti      & -   & 2   & -   & -   & -   & -   & 2                  \\ \hline
                Lapenna       & -   & -   & -   & -   & 6   & -   & 6                  \\ \hline
                Pan           & -   & -   & -   & -   & -   & 3   & 3                  \\ \hline
                Pinarello     & -   & -   & -   & 3   & -   & -   & 3                  \\ \hline
                Rizzolo       & 1,5 & -   & -   & -   & -   & -   & 1,5                \\ \hline
                Totale        & 1,5 & 2   & 2   & 3   & 6   & 6   &                    \\ \hline
                \end{tabular}
                \caption{Consuntivo orario per ruolo nello Sprint 18}
            \end{table}

            % TODO modificare colonna "Ore Rimanenti"
            \paragraph{Prospetto Economico}\mbox{}\\
            \begin{table}[H]
                \centering
                \begin{tabular}{|c|c|c|c|}
                \hline
                \textbf{Ruolo}  & \textbf{Ore}  & \textbf{Costo}      & \textbf{Ore rimanenti}   \\ \hline
                Responsabile    & 1,5           & €45                 & 5,5                      \\ \hline
                Amministratore  & 2             & €40                 & 6                        \\ \hline
                Analista        & 2             & €50                 & 5                        \\ \hline
                Progettista     & 3             & €75                 & 8,5                      \\ \hline
                Programmatore   & 6             & €90                 & 19,5                     \\ \hline
                Verificatore    & 6             & €90                 & 8                        \\ \hline
                \textbf{Totale} & \textbf{20,5} & \textbf{\euro340}   & \textbf{52,5}            \\ \hline
                \end{tabular}
                \caption{Prospetto economico e ore rimanenti}
            \end{table}
            
            \paragraph{Rischi effettivamente occorsi}
            \begin{itemize}
                \item RT1: Alcune difficoltà nell’integrazione di FAISS e LangChain hanno richiesto più tempo del previsto per essere risolte.
                \item RT2: Alcune parti del codice backend hanno richiesto chiarimenti e rifattorizzazioni per migliorare la leggibilità.
                \item RO1: Alcune attività sono state riorganizzate a causa di difficoltà impreviste, portando a una ripianificazione.
            \end{itemize}
            
            \paragraph{Retrospettiva}
            Lo sprint ha portato significativi miglioramenti al progetto, soprattutto in termini di funzionalità implementate e documentazione aggiornata. Tuttavia, sono emersi alcuni problemi di comunicazione e difficoltà tecniche che hanno richiesto iterazioni aggiuntive. Per migliorare nei prossimi sprint, si consiglia di:
            \begin{itemize}
                \item Definire meglio la comunicazione tra frontend e backend, stabilendo momenti fissi di allineamento.
                \item Migliorare la documentazione interna per facilitare l’adozione di nuove tecnologie.
                \item Pianificare con maggiore margine le attività che coinvolgono strumenti complessi come FAISS e LangChain.
            \end{itemize}
                
            %%%%%%%%%%%%%%%%%%%%%%%%%%%%%%%%%%%%%%%%%%%%%%%%%%%%%%%%%%%%%%
            \subsubsection{Sprint 19}
            \paragraph{Pianificazione}
                \begin{itemize}
                    \item inizio: 17 Marzo 2025
                    \item fine: 23 Marzo 2025
                \end{itemize}
        
            \paragraph{Obbiettivi dello Sprint}
            \begin{itemize}
                \item Completamento e testing dei controller, con una copertura del 85\%.
                \item Risoluzione di problemi nei test e nei repository, migliorando la stabilità del codice.
                \item Implementazione e correzione degli endpoint per la comunicazione tra frontend e backend.
                \item Sviluppo e rifinitura dei diagrammi UML, con particolare attenzione al diagramma esagonale.
                \item Introduzione della gestione della chat, compreso il salvataggio ed eliminazione delle conversazioni.
                \item Miglioramenti nella documentazione, con aggiornamenti alla specifica tecnica e la creazione di nuovi verbali.
            \end{itemize}
            
            \paragraph{Task Prioritari}
            I task principali affrontati durante lo sprint sono stati:
            \begin{itemize}
                \item Scrittura e verifica della sezione 3.2 della specifica tecnica.
                \item Creazione del verbale del 14 marzo 2025.
                \item Testing approfondito sui controller e servizi, con copertura oltre l'85\%.
                \item Correzioni ai diagrammi UML e miglioramenti visivi.
                \item Implementazione degli endpoint principali e correzione dei test di integrazione.
                \item Rifinitura del supporto per la gestione della chat.
                \item Revisione e miglioramento della connessione degli adapter.
            \end{itemize}
            
            \paragraph{Rischi attesi}
            \begin{itemize}
                \item RT1: Difficoltà nell’uso di nuove tecnologie, specialmente nell’integrazione dei test sugli adapter e nella gestione della chat.
                \item RT2: Sviluppo di codice poco comprensibile, soprattutto nei controller e nei repository.
                \item RC2: Problemi di comunicazione all’interno del gruppo, dovuti alla necessità di coordinare le modifiche tra frontend e backend.
                \item RO1: Incomprensioni nella pianificazione delle attività, con possibili ritardi nella gestione dei test e nella verifica della documentazione.
            \end{itemize}

            \paragraph{Preventivo}\mbox{}\\
            \begin{table}[H]
                \centering
                \begin{tabular}{|c|c|c|c|c|c|c|c|}
                \hline
                              & Re  & Am  & An  & Pj  & Pg  & Ve  & Totale per persona \\ \hline
                Cotti Cottini & -   & 1,5 & -   & -   & -   & -   & 1,5                \\ \hline
                Di Pietro     & -   & -   & -   & -   & -   & 3   & 3                  \\ \hline
                Diviesti      & -   & -   & -   & -   & 6,5 & -   & 6,5                \\ \hline
                Lapenna       & -   & -   & -   & -   & -   & 3   & 3                  \\ \hline
                Pan           & 2   & -   & -   & -   & -   & -   & 2                  \\ \hline
                Pinarello     & -   & -   & 1,5 & -   & -   & -   & 1,5                \\ \hline
                Rizzolo       & -   & -   & -   & 3   & -   & -   & 3                  \\ \hline
                Totale        & 2   & 1,5 & 1,5 & 3   & 6,5 & 6   &                    \\ \hline
                \end{tabular}
                \caption{Preventivo orario per ruolo nello Sprint 19}
            \end{table}


            \begin{figure}[H]
                \centering
                \begin{tikzpicture}
                    \pie[polar]{10/Responsabile, 7/Amministratore, 7/Analista, 15/Progettista, 32/Programmatore, 29/Verificatore}
                \end{tikzpicture}
                \caption{Ripartizione in percentuale dei ruoli nello Sprint 19}
            \end{figure}


            % TODO copiala dal preventivo (per ora li facciamo uguali)
            \paragraph{Consuntivo}\mbox{}\\
            \begin{table}[H]
                \centering
                \begin{tabular}{|c|c|c|c|c|c|c|c|}
                \hline
                              & Re  & Am  & An  & Pj  & Pg  & Ve  & Totale per persona \\ \hline
                Cotti Cottini & -   & 1,5 & -   & -   & -   & -   & 1,5                \\ \hline
                Di Pietro     & -   & -   & -   & -   & -   & 3   & 3                  \\ \hline
                Diviesti      & -   & -   & -   & -   & 6,5 & -   & 6,5                \\ \hline
                Lapenna       & -   & -   & -   & -   & -   & 3   & 3                  \\ \hline
                Pan           & 2   & -   & -   & -   & -   & -   & 2                  \\ \hline
                Pinarello     & -   & -   & 1,5 & -   & -   & -   & 1,5                \\ \hline
                Rizzolo       & -   & -   & -   & 3   & -   & -   & 3                  \\ \hline
                Totale        & 2   & 1,5 & 1,5 & 3   & 6,5 & 6   &                    \\ \hline
                \end{tabular}
                \caption{Consuntivo orario per ruolo nello Sprint 19}
            \end{table}

            % TODO modificare colonna "Ore Rimanenti"
            \paragraph{Prospetto Economico}\mbox{}\\
            \begin{table}[H]
                \centering
                \begin{tabular}{|c|c|c|c|}
                \hline
                \textbf{Ruolo}  & \textbf{Ore}  & \textbf{Costo}      & \textbf{Ore rimanenti}   \\ \hline
                Responsabile    & 2             & €60                 & 3,5                      \\ \hline
                Amministratore  & 1,5           & €30                 & 4,5                        \\ \hline
                Analista        & 1,5           & €37,5               & 3,5                      \\ \hline
                Progettista     & 3             & €75                 & 5,5                      \\ \hline
                Programmatore   & 6,5           & €97,5               & 13                       \\ \hline
                Verificatore    & 6             & €90                 & 12                        \\ \hline
                \textbf{Totale} & \textbf{20,5} & \textbf{\euro390}   & \textbf{42}            \\ \hline
                \end{tabular}
                \caption{Prospetto economico e ore rimanenti}
            \end{table}
            
            \paragraph{Rischi effettivamente occorsi}
            \begin{itemize}
                \item RT1: Alcune difficoltà tecniche sono emerse nella gestione degli adapter e nella loro integrazione con i test.
                \item RT2: Alcune parti del codice hanno richiesto refactoring per migliorarne la leggibilità e manutenibilità.
                \item RC2: Problemi di comunicazione tra i membri del team hanno reso necessario un maggiore allineamento sulle implementazioni degli endpoint.
                \item RO1: Alcuni task sono stati ridefiniti e riorganizzati a causa di difficoltà riscontrate nella gestione della chat e nei repository.
            \end{itemize}
            
            \paragraph{Retrospettiva}
            Lo sprint ha portato a significativi miglioramenti in termini di stabilità del codice, copertura dei test e qualità della documentazione. Tuttavia, alcuni problemi di comunicazione e difficoltà tecniche hanno causato ritardi su alcune funzionalità chiave. Per i prossimi sprint, il team ha identificato le seguenti azioni di miglioramento:
            \begin{itemize}
                \item Migliorare il coordinamento tra i membri del team, specialmente nelle fasi di testing e integrazione.
                \item Documentare meglio le implementazioni più complesse per ridurre il rischio di codice poco comprensibile.
                \item Definire momenti fissi di allineamento per evitare ritardi dovuti a incomprensioni nella pianificazione.
            \end{itemize}
                
            %%%%%%%%%%%%%%%%%%%%%%%%%%%%%%%%%%%%%%%%%%%%%%%%%%%%%%%%%%%%%%
            \subsubsection{Sprint 20}
            \paragraph{Pianificazione}
                \begin{itemize}
                    \item inizio: 24 Marzo 2025
                    \item fine: 30 Marzo 2025
                \end{itemize}
            
            \paragraph{Obbiettivi dello Sprint}
            \begin{itemize}
                \item Implementazione della gestione della sessione utente e funzionalità di logout.
                \item Rifinitura e testing della chat, inclusa l'eliminazione delle conversazioni e miglioramenti all’interfaccia.
                \item Introduzione del supporto a react-markdown e miglioramenti nella gestione degli errori API.
                \item Miglioramenti al database, con aggiornamenti dello schema per supportare i messaggi e l'assegnazione di utenti alle conversazioni.
                \item Revisione della documentazione, con aggiornamenti ai verbali e alla lista delle tecnologie utilizzate.
                \item Miglioramento del rating system con l'implementazione del feedback Thumbs Up/Down.
                \item Fix di test, con copertura aumentata al 74\%.
            \end{itemize}
            
            \paragraph{Task Prioritari}
            \begin{itemize}
                \item Sviluppo e testing della gestione delle conversazioni e dei messaggi nella chat.
                \item Implementazione e miglioramento della gestione delle sessioni utente.
                \item Rifinitura dell'architettura e aggiornamento della documentazione tecnica.
                \item Revisione e test delle funzionalità di assistenza.
                \item Refactoring del codice, con miglioramenti nei modelli di conversazione e nei servizi.
                \item Fix dei diagrammi UML e miglioramenti ai template grafici.
                \item Risoluzione di problemi nei test di integrazione e miglioramento della copertura dei test.
            \end{itemize}
            
            \paragraph{Rischi attesi}
            \begin{itemize}
                \item RT1: Difficoltà nell’uso di nuove tecnologie, in particolare nella gestione delle sessioni utente e nel miglioramento del chatbot.
                \item RT2: Sviluppo di codice poco comprensibile a causa delle numerose modifiche e refactoring.
                \item RC2: Problemi di comunicazione all’interno del gruppo durante l'integrazione delle nuove funzionalità.
            \end{itemize}

            \paragraph{Preventivo}\mbox{}\\
            \begin{table}[H]
                \centering
                \begin{tabular}{|c|c|c|c|c|c|c|c|}
                \hline
                              & Re  & Am  & An  & Pj  & Pg  & Ve  & Totale per persona \\ \hline
                Cotti Cottini & -   & -   & -   & 2,5 & -   & -   & 2,5                \\ \hline
                Di Pietro     & -   & 2   & -   & -   & -   & -   & 2                  \\ \hline
                Diviesti      & -   & -   & -   & -   & -   & 3   & 3                  \\ \hline
                Lapenna       & -   & -   & -   & -   & 6,5 & -   & 6,5                \\ \hline
                Pan           & -   & -   & -   & -   & -   & 3   & 3                  \\ \hline
                Pinarello     & 2   & -   & -   & -   & -   & -   & 2                  \\ \hline
                Rizzolo       & -   & -   & 2   & -   & -   & -   & 2                  \\ \hline
                Totale        & 2   & 2   & 2   & 2,5 & 6,5 & 6   &                    \\ \hline
                \end{tabular}
                \caption{Preventivo orario per ruolo nello Sprint 20}
            \end{table}


            \begin{figure}[H]
                \centering
                \begin{tikzpicture}
                    \pie[polar]{9/Responsabile, 10/Amministratore, 9/Analista, 12/Progettista, 31/Programmatore, 29/Verificatore}
                \end{tikzpicture}
                \caption{Ripartizione in percentuale dei ruoli nello Sprint 20}
            \end{figure}


            % TODO copiala dal preventivo (per ora li facciamo uguali)
            \paragraph{Consuntivo}\mbox{}\\
            \begin{table}[H]
                \centering
                \begin{tabular}{|c|c|c|c|c|c|c|c|}
                \hline
                                & Re  & Am  & An  & Pj  & Pg  & Ve  & Totale per persona \\ \hline
                Cotti Cottini & -   & -   & -   & 2,5 & -   & -   & 2,5                \\ \hline
                Di Pietro     & -   & 2   & -   & -   & -   & -   & 2                  \\ \hline
                Diviesti      & -   & -   & -   & -   & -   & 3   & 3                  \\ \hline
                Lapenna       & -   & -   & -   & -   & 6,5 & -   & 6,5                \\ \hline
                Pan           & -   & -   & -   & -   & -   & 3   & 3                  \\ \hline
                Pinarello     & 2   & -   & -   & -   & -   & -   & 2                  \\ \hline
                Rizzolo       & -   & -   & 2   & -   & -   & -   & 2                  \\ \hline
                Totale        & 2   & 2   & 2   & 2,5 & 6,5 & 6   &                    \\ \hline
                \end{tabular}
                \caption{Consuntivo orario per ruolo nello Sprint 20}
            \end{table}

            % TODO modificare colonna "Ore Rimanenti"
            \paragraph{Prospetto Economico}\mbox{}\\
            \begin{table}[H]
                \centering
                \begin{tabular}{|c|c|c|c|}
                \hline
                \textbf{Ruolo}  & \textbf{Ore}  & \textbf{Costo}      & \textbf{Ore rimanenti} \\ \hline
                Responsabile    & 2             & €60                 & 1,5                    \\ \hline
                Amministratore  & 2             & €40                 & 2,5                    \\ \hline
                Analista        & 2             & €50                 & 1,5                    \\ \hline
                Progettista     & 2,5           & €62,5               & 3                      \\ \hline
                Programmatore   & 6,5           & €97,5               & 6,5                    \\ \hline
                Verificatore    & 6             & €90                 & 6                      \\ \hline
                \textbf{Totale} & \textbf{21}   & \textbf{\euro400}   & \textbf{21}            \\ \hline
                \end{tabular}
                \caption{Prospetto economico e ore rimanenti}
            \end{table}
            
            \paragraph{Rischi effettivamente occorsi}
            \begin{itemize}
                \item RT1: Sono emerse difficoltà nell'integrazione del feedback utente e nella gestione delle sessioni.
                \item RT2: Alcune parti del codice sono risultate poco chiare e hanno richiesto rifattorizzazione per migliorare la leggibilità.
                \item RC2: Alcuni problemi di comunicazione hanno reso necessaria una maggiore sincronizzazione nelle modifiche al database e ai controller.
            \end{itemize}
            
            \paragraph{Retrospettiva}
            Lo sprint ha permesso di fare significativi progressi nelle funzionalità chiave del progetto, migliorando l'interfaccia, l'integrazione della chat e la gestione degli utenti. Tuttavia, ci sono stati alcuni ritardi e difficoltà tecniche che hanno rallentato l'implementazione di alcuni aspetti. Per i prossimi sprint, il team ha identificato le seguenti azioni di miglioramento:
            \begin{itemize}
                \item Rafforzare la comunicazione interna per migliorare il coordinamento tra le diverse aree di sviluppo.
                \item Documentare meglio le funzionalità critiche per facilitare la comprensione e la manutenzione del codice.
                \item Ottimizzare la gestione del tempo per ridurre l'impatto degli impegni esterni sulle scadenze del progetto.
                \item Introdurre più sessioni di testing intermedio per individuare problemi di integrazione in anticipo.
            \end{itemize}
                
            %%%%%%%%%%%%%%%%%%%%%%%%%%%%%%%%%%%%%%%%%%%%%%%%%%%%%%%%%%%%%%
            \subsubsection{Sprint 21}
            \paragraph{Pianificazione}
                \begin{itemize}
                    \item inizio: 31 Marzo 2025
                    \item fine: 6 Aprile 2025
                \end{itemize}
            
            \paragraph{Obbiettivi dello Sprint}
            \begin{itemize}
                \item Migliorare l'integrazione del chatbot con i template utente, completare il manuale utente e affinare il sistema di gestione delle metriche della dashboard.
            \end{itemize}
            
            
            \paragraph{Task Prioritari}  
            \begin{itemize}  
                \item Implementazione del recupero e gestione dei template utente nel chatbot  
                \item Completamento e revisione del manuale utente  
                \item Ottimizzazione dell'integrazione delle metriche nella dashboard  
                \item Refinement dell'interfaccia grafica (aggiornamento colori e layout chatbot)  
            \end{itemize}  
            
            \paragraph{Rischi attesi}  
            \begin{itemize}  
                \item RT1 Difficoltà nell’uso di nuove tecnologie  
                \item RO1 Incomprensioni nella pianificazione delle attività  
            \end{itemize}

            \paragraph{Preventivo}\mbox{}\\
            \begin{table}[H]
                \centering
                \begin{tabular}{|c|c|c|c|c|c|c|c|}
                \hline
                              & Re  & Am  & An  & Pj  & Pg  & Ve  & Totale per persona \\ \hline
                Cotti Cottini & -   & -   & 1,5 & -   & -   & -   & 1,5                \\ \hline
                Di Pietro     & -   & -   & -   & 3   & -   & -   & 3                  \\ \hline
                Diviesti      & -   & 2,5 & -   & -   & -   & -   & 2,5                \\ \hline
                Lapenna       & -   & -   & -   & -   & -   & 3   & 3                  \\ \hline
                Pan           & -   & -   & -   & -   & 6,5 & -   & 6,5                \\ \hline
                Pinarello     & -   & -   & -   & -   & -   & 3   & 3                  \\ \hline
                Rizzolo       & 1,5 & -   & -   & -   & -   & -   & 1,5                \\ \hline
                Totale        & 1,5 & 2,5 & 1,5 & 3   & 6,5 & 6   &                    \\ \hline
                \end{tabular}
                \caption{Preventivo orario per ruolo nello Sprint 21}
            \end{table}


            \begin{figure}[H]
                \centering
                \begin{tikzpicture}
                    \pie[polar]{7/Responsabile, 12/Amministratore, 7/Analista, 14/Progettista, 31/Programmatore, 29/Verificatore}
                \end{tikzpicture}
                \caption{Ripartizione in percentuale dei ruoli nello Sprint 21}
            \end{figure}


            % TODO copiala dal preventivo (per ora li facciamo uguali)
            \paragraph{Consuntivo}\mbox{}\\
            \begin{table}[H]
                \centering
                \begin{tabular}{|c|c|c|c|c|c|c|c|}
                \hline
                              & Re  & Am  & An  & Pj  & Pg  & Ve  & Totale per persona \\ \hline
                Cotti Cottini & -   & -   & 1,5 & -   & -   & -   & 1,5                \\ \hline
                Di Pietro     & -   & -   & -   & 3   & -   & -   & 3                  \\ \hline
                Diviesti      & -   & 2,5 & -   & -   & -   & -   & 2,5                \\ \hline
                Lapenna       & -   & -   & -   & -   & -   & 3   & 3                  \\ \hline
                Pan           & -   & -   & -   & -   & 6,5 & -   & 6,5                \\ \hline
                Pinarello     & -   & -   & -   & -   & -   & 3   & 3                  \\ \hline
                Rizzolo       & 1,5 & -   & -   & -   & -   & -   & 1,5                \\ \hline
                Totale        & 1,5 & 2,5 & 1,5 & 3   & 6,5 & 6   &                    \\ \hline
                \end{tabular}
                \caption{Consuntivo orario per ruolo nello Sprint 21}
            \end{table}

            % TODO modificare colonna "Ore Rimanenti"
            \paragraph{Prospetto Economico}\mbox{}\\
            \begin{table}[H]
                \centering
                \begin{tabular}{|c|c|c|c|}
                \hline
                \textbf{Ruolo}  & \textbf{Ore}  & \textbf{Costo}      & \textbf{Ore rimanenti} \\ \hline
                Responsabile    & 1,5           & €45                 & 0                      \\ \hline
                Amministratore  & 2,5           & €50                 & 0                      \\ \hline
                Analista        & 1,5           & €37,5               & 0                      \\ \hline
                Progettista     & 3             & €75                 & 0                      \\ \hline
                Programmatore   & 6,5           & €97,5               & 0                      \\ \hline
                Verificatore    & 6             & €90                 & 0                      \\ \hline
                \textbf{Totale} & \textbf{21}   & \textbf{\euro395}   & \textbf{0}            \\ \hline
                \end{tabular}
                \caption{Prospetto economico e ore rimanenti}
            \end{table}
            
            \paragraph{Rischi effettivamente occorsi}  
            \begin{itemize}  
                \item RT1 Difficoltà nell’uso di nuove tecnologie: alcune funzionalità del chatbot hanno richiesto più tempo del previsto a causa di problemi di integrazione con i template utente.  
                \item RO1 Incomprensioni nella pianificazione delle attività: alcune attività hanno subito ritardi per via della necessità di rivedere il manuale utente e correggere alcune sezioni della documentazione.  
            \end{itemize}  
            
            \paragraph{Retrospettiva}  
            Nonostante alcune difficoltà iniziali nell'uso delle nuove tecnologie e nella gestione dei tempi, il team è riuscito a portare a termine le attività prioritarie. La documentazione è stata migliorata e il chatbot ha ricevuto aggiornamenti significativi.
                
            %%%%%%%%%%%%%%%%%%%%%%%%%%%%%%%%%%%%%%%%%%%%%%%%%%%%%%%%%%%%%%

%%%%%%%%%%%%%%%%%%%%%%%%%%%%%%%%%%%%%%%%%%%%%%%%%%%%%%%%%%%%%%%%%%%%%%%%%%%%%%%%%%%%%%

% \paragraph{Preventivo}
% \paragraph{Consuntivo}
% \paragraph{Prospetto Economico}

\end{document}
